% Chapter 9 --- Beyond the paper: a call-by-value model (Melliès §7.4).
%
% This chapter documents the call-by-value (CBV) development, which goes
% BEYOND the Ehrhard–Geoffroy paper: their §9 is the LNL/Seely model, and
% CBV is the paper's STATED FUTURE WORK (conclusion: "In future work we
% will explain how this model can be used for interpreting call-by-value
% or even call-by-push-value … languages …").  The formalisation pushes in
% that direction along the Eilenberg–Moore / coalgebra route of
% Melliès, *Categorical Semantics of Linear Logic*, §7.4 (EM/coalgebra =
% CBV; co-Kleisli = CBN).
%
% Mirrors:
%   theories/homs/em_cat.v             EM(!) (ICones_EM), the cofree functor
%                                      !̃ = bang_cofree, the forgetful U, and
%                                      the adjunction U ⊣ !̃ (adj_phi/adj_psi,
%                                      round-trips, naturality, triangles)
%   theories/homs/em_seely_comonoid.v  the Seely comonoid d_bang/e_bang on !A,
%                                      LC2 (comonoid laws), LC3 (coalgebra-
%                                      morphism), LC4 (dig a comonoid morphism);
%                                      tens_cofree / unit_cofree
%   theories/homs/em_cartesian.v       the FULL EM(!) is cartesian with product
%                                      carried by the LINEAR ⊗ (not &);
%                                      EMComon_all unconditional via the
%                                      Melliès Prop 26-28 / Cor 17/20
%                                      retraction-and-lifting argument
%                                      (diagram81, coalg_mor_lift,
%                                      coalg_d_is_mor_gen);
%                                      ICones_EM_cartesian, EM_prod/EM_term,
%                                      m_bang; em_pair_mor / em_proj1_mor /
%                                      em_proj2_mor / em_term_mor; the
%                                      transported commutative-comonoid pair
%                                      (coalg_d, coalg_e) carried by every
%                                      !-coalgebra (Melliès Prop 28 / Cor 20).
%   theories/programs/infra/cbv_adjunction.v   the (lax symmetric) monoidal
%                                      adjunction U ⊣ !̃ between ICones (linear)
%                                      and the FULL category EM(!) of
%                                      !-coalgebras (Melliès Prop 29);
%                                      em_pair_mor_proj_id (cartesian η at
%                                      the icones level, Fox 1976 / Prop 28);
%                                      bang_m / bang_e0 / lax coherence /
%                                      monoidal counit laws; the CBV_Model
%                                      record + the witness ICones_CBV.
%   theories/homs/fmeas_lax.v          the FMeas lax symmetric monoidal map
%                                      fmeas_lax X Y : FMeas X ⊗ FMeas Y →
%                                      FMeas (X × Y) bundled as an icones_hom
%                                      via tensor_uncurry of the bilinear lift;
%                                      Dirac identity fmeas_lax_dirac; uses
%                                      bilin.v's int_to_linhom_pres_path_in_cone
%                                      (the now-discharged path-preservation
%                                      follow-up to paper Thm 6.1)
%   theories/programs/infra/bool_cone.v   the 2-point sub-probability cone
%                                      (paper §4.4 / Thm 4.24 coproduct
%                                      cone_one ⊕ cone_one), bool_case
%                                      universal co-pairing; the boolean
%                                      cascade infrastructure used by ne_if.
%   theories/programs/infra/bool_case_hom.v   bool_case at the linhom level
%                                      (bool_case_linhom), unit-ball-free
%                                      variant, alpha/beta decomposition.
%   theories/programs/ppl.v            a higher-order, DIRECT-STYLE
%                                      (Plotkin/Girard CBV), SINGLE-SORT,
%                                      MULTI-VARIABLE, NAMED-VARIABLE
%                                      QBS-MIRROR PPL: one inductive
%                                      [named_expr Γ τ] with constructors
%                                      ne_var / ne_tt / ne_pair / ne_fst /
%                                      ne_snd / ne_lam / ne_app / ne_let /
%                                      ne_real / ne_add / ne_mul /
%                                      ne_sample / ne_score / ne_true /
%                                      ne_false / ne_bernoulli / ne_if /
%                                      ne_fix / ne_fix_mr.  Types
%                                      tunit / tbase X / tprod / tfun /
%                                      tbool; no tprob marker, no Ret, no
%                                      ne_bind in the source.  This file
%                                      now hosts ONLY the shared surface
%                                      syntax + cross-cutting helpers
%                                      (add_meas, mul_meas, add_lift,
%                                      mul_lift, score_lift, const_icones,
%                                      bernoulli, etc.) used by BOTH the
%                                      CBV interpreter of ppl_cbv.v and
%                                      the CBN interpreter of ppl_cbn*.v.
%   theories/programs/ppl_cbv.v        the CLEAN CBV interpreter — the
%                                      linhom-valued, comonoid-primitive
%                                      interpretation of named_expr.  Every
%                                      expression denotes a linear morphism
%                                      [eD M : linhom_car Γ̂ τ̂] in ICones,
%                                      NOT a coalgebra morphism.  The
%                                      coalgebra structure of the context
%                                      is used only through the transported
%                                      commutative-comonoid pair
%                                      (coalg_d, coalg_e) carried by every
%                                      !-coalgebra (Melliès Prop 28 /
%                                      Cor 20).  No Moggi monad in the
%                                      interpretation: no Tobj, no
%                                      tunit_eta, no kbind / kcomp.
%                                      tyD (tfun A B) = !̃(U⟦A⟧ ⊸ U⟦B⟧),
%                                      a SINGLE !̃ on the function type
%                                      (NO Tobj wrap on the codomain).
%                                      The two boundaries with !̃ are ne_lam
%                                      (adj_psi of tensor_curry, i.e. the
%                                      coalg_str of the context's coalgebra)
%                                      and ne_app (der to extract the
%                                      function value out of its !̃-wrap).
%                                      ne_fix / ne_fix_mr now route through
%                                      Yfix_fun_lin of em_fix.v — the CBV
%                                      Kleene value-fixpoint on the clean
%                                      linhom cone (no Tobj / !̃U wrap on B).
%   theories/programs/infra/em_fix.v   the CBV Kleene value-fixpoint on
%                                      the clean linhom cone: Phi_fun /
%                                      Yfix_fun_lin and their ball / monot.
%                                      / ω-continuity laws; parametric in
%                                      any iCone codomain B and any
%                                      diagonal δ on the context iCone,
%                                      so applies to bang_cofree / FMeas_
%                                      coalgebra / bool_cone_coalg / EM_prod
%                                      uniformly.
%   theories/programs/infra/bool_cone_coalg.v   the §9.7-style 2-point
%                                      !-coalgebra on bool_cone_car (a hand
%                                      built specialisation of FMeas_
%                                      coalgebra to the boolean carrier,
%                                      without exposing a measurableType R
%                                      on Ar): bool_coalg_str, the two
%                                      coalgebra laws (counit / coassoc),
%                                      and the universal-property dispatch
%                                      lemma bool_cone_dispatch.
%   theories/programs/examples.v       eleven surface PPL programs (no
%                                      denotation reduction lemmas; this
%                                      file is now the shared surface
%                                      program pool consumed by both the
%                                      CBV and the CBN interpreters):
%                                      QBS headlines ex_random_constant,
%                                      ex_random_linear, ex_bayes_linear;
%                                      recursive examples ex_loop, ex_geom,
%                                      ex_almost_loop; boolean sanity
%                                      checks ex_true, ex_false,
%                                      ex_fair_coin, ex_if_demo.
%
% Status note (post-cleanup): every section of this chapter is AXIOM-FREE
% (only the three classical boolp axioms): the Eilenberg–Moore category
% and cofree adjunction; the Seely comonoid LC2–LC4; the FULL cartesian
% EM(!) carried by the linear ⊗ (EMComon_all unconditional, Cor 20); the
% cartesian η-rule em_pair_mor_proj_id (Fox 1976 / Prop 28 at the icones
% level); the (lax symmetric) monoidal adjunction ICones_CBV; the FMeas
% lax monoidal map and its Dirac identity; the boolean cascade
% (bool_cone_car / bool_case / bool_case_linhom).  The interpreter
% itself, ppl_cbv.v, defines eD as a LINHOM-VALUED, comonoid-primitive
% recursion on named_expr; it is axiom-free, with ne_fix / ne_fix_mr
% routed through Yfix_fun_lin of em_fix.v (the Kleene value-fixpoint on
% the clean linhom cone) and tbool routed through bool_cone_coalg of
% bool_cone_coalg.v (the §9.7-style 2-point !-coalgebra).  The previous
% Moggi-monad presentation in the now-deleted cbv.v / earlier em_fix.v /
% em_fix_arr.v / case_em_red.v / curry_kbind.v / cbv_outer_pt.v /
% em_continuity.v / ex_geom_step.v / ex_loop_arr.v / ex_almost_loop_arr.v /
% ex_almost_loop_step.v / ex_headlines.v / ex_random_dist.v /
% ppl_cbv_geom_dist.v / ex_even_odd.v has been retired; the linhom-valued
% reading on the comonoid-primitive recursion is lighter to formalise and
% equivalent at the level of denotation (the cofree !̃ only surfaces at
% function arrows, exactly per the linear-map view of programs).

%%%%%%%%%%%%%%%%%%%%%%%%%%%%%%%%%%%%%%%%%%%%%%%%%%%%%%%%%%%%%%%%%%%%%%%%
\chapter{Beyond the paper: a call-by-value model (Melli\`es \S7.4)}%
\label{ch:cbv}

\paragraph{This chapter is beyond the Ehrhard--Geoffroy paper.}
Their \S9 is the LNL/Seely model of intuitionistic linear logic
(Chapter~\ref{ch:exponential}); a \emph{call-by-value} reading is the
paper's explicitly \emph{stated future work} --- the conclusion writes
that ``in future work we will explain how this model can be used for
interpreting call-by-value or even call-by-push-value $\ldots$
languages.'' The development of this chapter takes a step in that
direction, following Melli\`es' \emph{Categorical Semantics of Linear
Logic}~\cite{mellies-categorical-semantics}, \S7.4: the
\emph{Eilenberg--Moore} (coalgebra) route to the exponential gives a
\emph{call-by-value} model, exactly as the co-Kleisli route gives
call-by-name (the latter is the cartesian closed $\mathbf{SCones}$ of
Chapter~\ref{ch:stable}). Nothing here is filed as a paper-\S\ result; it
is a \emph{beyond-the-paper} construction on top of the axiom-free Seely
category of \autoref{thm:seely}.

\paragraph{What is, and is not, done.}
The CBV \emph{model} is built and \textbf{axiom-free} (only the three
classical \texttt{boolp} axioms): the \emph{full} cartesian
Eilenberg--Moore category, the cartesian $\eta$-rule at the icones level
(Fox 1976 / Melli\`es Proposition 28), and the (lax symmetric) monoidal
adjunction \rocq{Icones.programs.infra.cbv_adjunction.ICones_CBV}. The
cartesianness of $\mathrm{EM}(\mathord{!})$ is \emph{unconditional} ---
the predicate \rocq{Icones.homs.em_cartesian.EMComon} holds on every
$\mathord{!}$-coalgebra (\rocq{Icones.homs.em_cartesian.EMComon_all}),
following Melli\`es' Proposition 26--28 / Corollary 17/20 by the
structural retraction-and-lifting argument (\emph{not} the naïve
promoted-point reduction, which only computes on $x!$ and stalls on a
general carrier); $U \dashv \widetilde{\mathord{!}}$ is therefore a
\emph{linear / non-linear} adjunction with \emph{the} category of
$\mathord{!}$-coalgebras as the cartesian non-linear / value side. The
\emph{interpreter} of the higher-order PPL is
\texttt{theories/programs/ppl\_cbv.v}, a linhom-valued, comonoid-primitive
denotation on the shared surface inductive \texttt{named\_expr} of
\texttt{theories/programs/ppl.v}.

\paragraph{Linhom-valued, comonoid-primitive interpretation.}
The interpreter assigns to a term
$\Gamma \vdash M : \tau$ a \emph{linear} morphism in $\ICones$,
\[
  \sem{M}\;\;\in\;\;
  \mathrm{linhom}\bigl(U\sem{\Gamma},\,U\sem{\tau}\bigr),
\]
where $\sem{\Gamma}$ and $\sem{\tau}$ are $\mathord{!}$-coalgebras. The
denotation is \emph{not} a coalgebra morphism: a probabilistic program
need not commute with the comonoid duplicability structure carried by
its source coalgebra. The canonical witness is the let-then-pair
fragment $\,\mathtt{let}\,x \,{=}\, \mathtt{sample}\,\mu\,\mathtt{in}\,(x, x)\,$: the
share-then-measure denotation is a pushforward along the diagonal,
whereas the measure-then-pair denotation is the cartesian product
measure; the two distributions agree iff $\mu$ is a Dirac. Programs are
therefore modelled as linear maps and never asked to be coalgebra
morphisms; what every context coalgebra \emph{does} supply uniformly,
through Melli\`es' Proposition 28 / Corollary 20
(\rocq{Icones.homs.em_cartesian.EMComon_all}), is a commutative-comonoid
pair $(\delta_\Gamma, \varepsilon_\Gamma) =
(\mathtt{coalg\_d}, \mathtt{coalg\_e})$ on the underlying iCone, and
that pair is the only structural ingredient the interpreter ever uses on
the context side.

The cofree exponential $\widetilde{\mathord{!}}$ now surfaces at
\emph{exactly two} places, both at the function-type interpretation:
\begin{itemize}
  \item \rocq{Icones.programs.ppl.ne_lam} wraps a curried body
    $\mathrm{linhom}(U\sem{\Gamma} \otimes U\sem{A}, U\sem{B})$ ---
    obtained by SMCC currying --- into the cofree $\widetilde{\mathord{!}}$
    using the cofree adjunction's $\mathtt{adj\_psi}$ applied at the
    coalgebra structure of $\sem{\Gamma}$ (caller-side duplicability of
    function values);
  \item \rocq{Icones.programs.ppl.ne_app} extracts a function value from
    its $\widetilde{\mathord{!}}$-wrap via the counit
    $\mathrm{der} : \widetilde{\mathord{!}}L \to L$ on the function-cone
    $L = U\sem{A} \multimap U\sem{B}$ before tensor-uncurrying into
    evaluation.
\end{itemize}
Every other compositional clause uses only the SMC primitives (tensor,
braid, unitors, identity, composition) and the comonoid pair
$(\delta_\Gamma, \varepsilon_\Gamma)$. No Moggi monad
$T = \widetilde{\mathord{!}} \circ U$ ever appears in the codomain of a
type or in the interpretation: there is no $\mathtt{Tobj}$, no
$\mathtt{tunit\_eta}$, no Kleisli composition. The function type
\[
  \sem{\mathtt{tfun}\,A\,B}
    \;=\;
    \widetilde{\mathord{!}}\bigl(U\sem{A} \multimap U\sem{B}\bigr)
\]
carries a \emph{single} cofree $\widetilde{\mathord{!}}$ on the outside;
the codomain is bare.

\paragraph{The shared surface syntax.}
The surface inductive \rocq{Icones.programs.ppl.named_expr} is a
\emph{direct-style} (no value/computation term split, no \texttt{tprob}
type marker, no syntactic \texttt{return}, no syntactic \texttt{bind}),
\emph{single-sort}, \emph{multi-variable}, \emph{named-variable}
QBS-mirror calculus --- the modern PPL surface in the style of
Saito and Affeldt's APLAS 2023 named-variable
embedding\footnote{Saito and Affeldt, \emph{Towards a practical embedding
of a probabilistic programming language in Coq}, APLAS 2023.}. It is
defined once in \texttt{theories/programs/ppl.v} and consumed by
\emph{two} interpreters: the linhom-valued CBV interpreter of
\texttt{theories/programs/ppl\_cbv.v} described in this chapter, and
the SCones-valued CBN interpreter of \texttt{theories/programs/ppl\_cbn*.v}
of \autoref{ch:cbn}. The shared surface programs
(\texttt{theories/programs/examples.v}) include
\rocq{Icones.programs.examples.ex_random_constant} (the QBS
distribution-over-a-function-space example),
\rocq{Icones.programs.examples.ex_random_linear} (a distribution over
$\lambda x.\,m\cdot x + b$ for $m, b \sim \mu$), the
prior/score/observe Bayesian example
\rocq{Icones.programs.examples.ex_bayes_linear}, and the recursive
\rocq{Icones.programs.examples.ex_loop} / \texttt{ex\_geom} /
\texttt{ex\_almost\_loop}.

\paragraph{Where the material lives.}
\texttt{theories/homs/em\_cat.v} (the EM category, the cofree functor and
the cofree/forgetful adjunction), \texttt{theories/homs/em\_seely\_comonoid.v}
(the Seely comonoid on $\mathord{!}A$ and Melli\`es' conditions
LC2--LC4), \texttt{theories/homs/em\_cartesian.v} (the \emph{full}
cartesian $\mathrm{EM}(\mathord{!})$ carried by the linear $\otimes$,
with the unconditional \rocq{Icones.homs.em_cartesian.EMComon_all} and the
transported comonoid pair \rocq{Icones.homs.em_cartesian.coalg_d} /
\rocq{Icones.homs.em_cartesian.coalg_e}),
\texttt{theories/programs/infra/cbv\_adjunction.v} (the LNL monoidal
adjunction, the cartesian $\eta$-rule
\rocq{Icones.programs.infra.cbv_adjunction.em_pair_mor_proj_id}, and the
\texttt{CBV\_Model} bundle), \texttt{theories/homs/fmeas\_lax.v} (the
FMeas lax symmetric monoidal map at the \texttt{icones\_hom} level,
built on the discharged
\rocq{Icones.homs.bilin.int_to_linhom_pres_path_in_cone}),
\texttt{theories/programs/infra/bool\_cone.v} and
\texttt{theories/programs/infra/bool\_case\_hom.v} (the 2-point boolean
cone and its universal co-pairing, used by the $\mathtt{ne\_if}$ clause),
\texttt{theories/programs/ppl.v} (the shared higher-order direct-style
multi-variable named-variable surface syntax),
\texttt{theories/programs/ppl\_cbv.v} (the linhom-valued
comonoid-primitive CBV interpreter), and
\texttt{theories/programs/examples.v} (the eleven shared surface PPL
programs consumed by both interpreters).

%%%%%%%%%%%%%%%%%%%%%%%%%%%%%%%%%%%%%%%%%%%%%%%%%%%%%%%%%%%%%%%%%%%%%%%%
\section{The Eilenberg--Moore category and the cofree adjunction}%
\label{sec:cbv-em}

The exponential comonad $(\mathord{!}, \mathrm{der}, \mathrm{dig})$ of
\autoref{thm:comonad} has an Eilenberg--Moore category
$\mathrm{EM}(\mathord{!})$: its objects are $\mathord{!}$-coalgebras
$(A, a : A \to \mathord{!}A)$ (the minimal coalgebra layer of
\texttt{coalgebra.v}, paper Theorem 9.7), and its morphisms are coalgebra
morphisms. This section packages that category and the cofree/forgetful
adjunction; everything is a diagram chase from the comonad and coalgebra
laws --- no new analysis.

\begin{definition}[The Eilenberg--Moore category $\mathrm{EM}(\mathord{!})$]%
\label{def:cbv-em-cat}
\rocq{Icones.homs.em_cat.coalg_hom}\rocqok
\rocq{Icones.homs.em_cat.EM_Cat}\rocqok
\rocq{Icones.homs.em_cat.ICones_EM}\rocqok
\uses{thm:comonad, def:comonad-record}
A bundled coalgebra morphism \rocq{Icones.homs.em_cat.coalg_hom} is an
$\ICones$-morphism together with the coalgebra-morphism side-condition.
The identity and composition (\rocq{Icones.homs.em_cat.coalg_id},
\rocq{Icones.homs.em_cat.coalg_comp}) satisfy the category laws
(\rocq{Icones.homs.em_cat.coalg_compIl},
\rocq{Icones.homs.em_cat.coalg_compIr},
\rocq{Icones.homs.em_cat.coalg_compA}), bundled in the record
\rocq{Icones.homs.em_cat.EM_Cat} with canonical witness
\rocq{Icones.homs.em_cat.ICones_EM} --- mirroring the
\rocq{Icones.homs.bang.Comonad} record-with-witness convention.
\end{definition}

\begin{definition}[The cofree functor $\widetilde{\mathord{!}}$
and the forgetful functor $U$]\label{def:cbv-cofree}
\rocq{Icones.homs.em_cat.bang_cofree}\rocqok
\rocq{Icones.homs.em_cat.U_obj}\rocqok
\uses{def:cbv-em-cat, def:struct-maps}
The cofree coalgebra on $B$ is
$\widetilde{\mathord{!}}B = (\mathord{!}B, \mathrm{dig}_B)$
(\rocq{Icones.homs.em_cat.bang_cofree}), a coalgebra by the counit /
coassociativity laws \autoref{thm:comonad}; its action on a morphism is
$\widetilde{\mathord{!}}f = \mathord{!}f$
(\rocq{Icones.homs.em_cat.bang_cofree_hom}, a coalgebra morphism by
naturality of $\mathrm{dig}$), with functoriality
\rocq{Icones.homs.em_cat.bang_cofree_fmap_id} /
\rocq{Icones.homs.em_cat.bang_cofree_fmap_comp}. The forgetful functor
$U : \mathrm{EM}(\mathord{!}) \to \ICones$ takes $(A, a)$ to its carrier
$A$ (\rocq{Icones.homs.em_cat.U_obj}, \rocq{Icones.homs.em_cat.U_mor}).
\end{definition}

\begin{theorem}[The cofree adjunction
$U \dashv \widetilde{\mathord{!}}$]\label{thm:cbv-cofree-adj}
\rocq{Icones.homs.em_cat.adj_phi}\rocqok
\rocq{Icones.homs.em_cat.adj_psi}\rocqok
\rocq{Icones.homs.em_cat.adj_triangleL}\rocqok
\rocq{Icones.homs.em_cat.adj_triangleR}\rocqok
\uses{def:cbv-cofree}
The forgetful functor is left adjoint to the cofree functor,
$U \dashv \widetilde{\mathord{!}}$, presented by the natural
bijection of hom-sets
\[
  \mathrm{EM}(\mathord{!})\bigl(\gamma,\ \widetilde{\mathord{!}}B\bigr)
  \;\cong\;
  \ICones\bigl(U\gamma,\ B\bigr),
\]
namely \rocq{Icones.homs.em_cat.adj_phi} and
\rocq{Icones.homs.em_cat.adj_psi}, mutually inverse
(\rocq{Icones.homs.em_cat.adj_phiK},
\rocq{Icones.homs.em_cat.adj_psiK}) and natural in both arguments
(\rocq{Icones.homs.em_cat.adj_phi_natL},
\rocq{Icones.homs.em_cat.adj_phi_natR}). The unit is the coalgebra
structure map (\rocq{Icones.homs.em_cat.adj_unit}) and the counit is the
dereliction $\mathrm{der}$ (\rocq{Icones.homs.em_cat.adj_counit}); the
triangle identities are \rocq{Icones.homs.em_cat.adj_triangleL} and
\rocq{Icones.homs.em_cat.adj_triangleR}.
\end{theorem}

%%%%%%%%%%%%%%%%%%%%%%%%%%%%%%%%%%%%%%%%%%%%%%%%%%%%%%%%%%%%%%%%%%%%%%%%
\section{The Seely comonoid on \texorpdfstring{$\mathord{!}A$}{!A} (LC2--LC4)}\label{sec:cbv-comonoid}

To organise the cartesian structure we first equip every cofree object
$\mathord{!}A$ with the commutative-comonoid structure of a Melli\`es
\emph{linear category} (conditions LC2--LC4 of
\cite{mellies-categorical-semantics}, \S7.4), transported from the
cartesian diagonal / terminal of $(\mathbin{\&}, \top)$ through the Seely
isomorphisms $\mathrm{Seely2}$ / $\mathrm{Seely0}$ of
\autoref{def:seely2}. As throughout \autoref{ch:exponential}, the laws
are diagram chases on promoted points, reduced by the
$\texttt{bang\_ext}$ / $\texttt{tens\_excl\_charact}$ extensionality and
the $\mathrm{Seely2E}$ / $\mathrm{Seely0E}$ computation laws.

\begin{definition}[The comultiplication $d$ and counit $e$ on
$\mathord{!}A$]\label{def:cbv-d-e}
\rocq{Icones.homs.em_seely_comonoid.d_bang}\rocqok
\rocq{Icones.homs.em_seely_comonoid.e_bang}\rocqok
\uses{def:seely2, def:cbv-cofree}
The comultiplication
$d_A : \mathord{!}A \lin \mathord{!}A \otimes \mathord{!}A$ is
$\mathrm{Seely2}^{-1} \circ \mathord{!}(\Delta_A)$ for the cartesian
diagonal $\Delta_A$ (\rocq{Icones.homs.em_seely_comonoid.diag}), and the
counit $e_A : \mathord{!}A \lin 1$ is
$\mathrm{Seely0}^{-1} \circ \mathord{!}(\langle\rangle_A)$ for the
terminal map (\rocq{Icones.homs.em_seely_comonoid.term}). Their
computation laws on a promoted point $x!$ (for $\cnorm{x} \leq 1$) are
$d_A(x!) = x! \otimes x!$ (\rocq{Icones.homs.em_seely_comonoid.d_bang_prom})
and $e_A(x!) = 1$ (\rocq{Icones.homs.em_seely_comonoid.e_bang_prom}).
\end{definition}

\begin{lemma}[LC2: $(\mathord{!}A, d_A, e_A)$ is a commutative
comonoid]\label{lem:cbv-lc2}
\rocq{Icones.homs.em_seely_comonoid.comonoid_coassoc}\rocqok
\rocq{Icones.homs.em_seely_comonoid.comonoid_counitL}\rocqok
\rocq{Icones.homs.em_seely_comonoid.comonoid_counitR}\rocqok
\rocq{Icones.homs.em_seely_comonoid.comonoid_cocomm}\rocqok
\uses{def:cbv-d-e, lem:tens-excl-charact}
The triple satisfies coassociativity
(\rocq{Icones.homs.em_seely_comonoid.comonoid_coassoc}), the two counit
laws (\rocq{Icones.homs.em_seely_comonoid.comonoid_counitL},
\rocq{Icones.homs.em_seely_comonoid.comonoid_counitR}) and
cocommutativity (\rocq{Icones.homs.em_seely_comonoid.comonoid_cocomm}).
\end{lemma}

\begin{proof}
Each law is reduced, by \autoref{lem:tens-excl-charact}, to agreement on
a promoted pure tensor; both sides then compute by
\rocq{Icones.homs.em_seely_comonoid.d_bang_prom} /
\rocq{Icones.homs.em_seely_comonoid.e_bang_prom} and the structural-iso
laws of the tensor.
\end{proof}

\begin{lemma}[LC3 / LC4: $d$, $e$ are coalgebra morphisms and
$\mathrm{dig}$ is a comonoid morphism]\label{lem:cbv-lc34}
\rocq{Icones.homs.em_seely_comonoid.d_bang_is_coalg_mor}\rocqok
\rocq{Icones.homs.em_seely_comonoid.e_bang_is_coalg_mor}\rocqok
\rocq{Icones.homs.em_seely_comonoid.dig_comonoid_mult}\rocqok
\rocq{Icones.homs.em_seely_comonoid.dig_comonoid_counit}\rocqok
\uses{lem:cbv-lc2}
LC3: $d_A$ and $e_A$ are morphisms of coalgebras
(\rocq{Icones.homs.em_seely_comonoid.d_bang_is_coalg_mor},
\rocq{Icones.homs.em_seely_comonoid.e_bang_is_coalg_mor}), with the
target coalgebras the transported $\mathord{!}A \otimes \mathord{!}A$
(\rocq{Icones.homs.em_seely_comonoid.tens_cofree}) and $1$
(\rocq{Icones.homs.em_seely_comonoid.unit_cofree}). LC4: the
comultiplication $\mathrm{dig}_A$ is a comonoid morphism, compatible with
$d$ (\rocq{Icones.homs.em_seely_comonoid.dig_comonoid_mult}) and with $e$
(\rocq{Icones.homs.em_seely_comonoid.dig_comonoid_counit}).
\end{lemma}

\begin{remark}[The symmetric-monoidal level only]
The objects $\mathord{!}A \otimes \mathord{!}A$
(\rocq{Icones.homs.em_seely_comonoid.tens_cofree}) and $1$
(\rocq{Icones.homs.em_seely_comonoid.unit_cofree}) carry coalgebra
structures inherited across $\mathrm{Seely2}$ / $\mathrm{Seely0}$ --- the
symmetric-monoidal (LC1-level) structure of $\mathrm{EM}(\mathord{!})$.
This is \emph{not yet} the cartesian product, which is the subject of the
next section.
\end{remark}

%%%%%%%%%%%%%%%%%%%%%%%%%%%%%%%%%%%%%%%%%%%%%%%%%%%%%%%%%%%%%%%%%%%%%%%%
\section{The cartesian \texorpdfstring{$\mathrm{EM}(\mathord{!})$}{EM(!)} via the linear \texorpdfstring{$\otimes$}{tensor}}\label{sec:cbv-cartesian}

This section equips the \emph{full} Eilenberg--Moore category
$\mathrm{EM}(\mathord{!})$ with a cartesian structure whose binary
product is carried by the \emph{linear} tensor $\otimes$ (\textbf{not} the
cartesian $\mathbin{\&}$) and whose terminal object is the tensor unit
$1$. This is the hallmark of the EM/CBV reading: in the co-Kleisli (CBN)
category the product is $\mathbin{\&}$, while in the Eilenberg--Moore
(CBV) category the product is $\otimes$. The non-trivial point is that
this works on \emph{every} $\mathord{!}$-coalgebra, not only on a rich
sub-class --- Melli\`es' Proposition 26--28 and Corollary 17/20 ---
proved here by his structural retraction-and-lifting argument.

\begin{definition}[The lax-monoidal map, the product and terminal
coalgebras]\label{def:cbv-prod}
\rocq{Icones.homs.em_cartesian.m_bang}\rocqok
\rocq{Icones.homs.em_cartesian.EM_prod}\rocqok
\rocq{Icones.homs.em_cartesian.EM_term}\rocqok
\uses{thm:cbv-cofree-adj, lem:cbv-lc34, lem:tens-excl-charact}
The lax symmetric-monoidal structure map of $\mathord{!}$ is
$m_{A,B} : \mathord{!}A \otimes \mathord{!}B \lin \mathord{!}(A \otimes B)$
(\rocq{Icones.homs.em_cartesian.m_bang}), with
$m(x! \otimes y!) = (x \otimes y)!$
(\rocq{Icones.homs.em_cartesian.m_bang_prom}); it is a coalgebra morphism
(\rocq{Icones.homs.em_cartesian.m_bang_is_coalg_mor}). The product
coalgebra is $(A \otimes B,\ m \circ (a \otimes b))$
(\rocq{Icones.homs.em_cartesian.EM_prod}) and the terminal coalgebra is
$(1, \cdot)$ (\rocq{Icones.homs.em_cartesian.EM_term}, the
$\texttt{unit\_cofree}$ coalgebra). The Eq-88 transported comonoid maps
on \emph{every} coalgebra are
\rocq{Icones.homs.em_cartesian.coalg_d} /
\rocq{Icones.homs.em_cartesian.coalg_e}; these definitions are general,
but their laws are per-object (next).
\end{definition}

\begin{definition}[The comonoidality predicate $\mathtt{EMComon}$]%
\label{def:cbv-comon}
\rocq{Icones.homs.em_cartesian.EMComon}\rocqok
\uses{def:cbv-prod}
A coalgebra $P$ is \emph{comonoidal} ---
\rocq{Icones.homs.em_cartesian.EMComon} holds of it --- when the
transported maps \rocq{Icones.homs.em_cartesian.coalg_d} /
\rocq{Icones.homs.em_cartesian.coalg_e} make its carrier a commutative
comonoid \emph{and} are coalgebra morphisms (Melli\`es' Propositions 26 /
27 and Corollary 20). This is the per-object property the cartesian
construction needs --- and it turns out to hold on every coalgebra
(\autoref{lem:cbv-emcomon-all}).
\end{definition}

\begin{theorem}[$\mathtt{EMComon}$ is unconditional --- Melli\`es
Proposition 26--28]\label{lem:cbv-emcomon-all}
\rocq{Icones.homs.em_cartesian.EMComon_all}\rocqok
\uses{def:cbv-comon, lem:cbv-lc2, lem:cbv-lc34}
For \emph{every} $\mathord{!}$-coalgebra $P = (A, a)$, the transported
maps \rocq{Icones.homs.em_cartesian.coalg_d} /
\rocq{Icones.homs.em_cartesian.coalg_e} satisfy all the comonoid /
coalgebra-morphism laws --- i.e.\ \rocq{Icones.homs.em_cartesian.EMComon_all}.
The cofree coalgebras $\widetilde{\mathord{!}}B$
(\rocq{Icones.homs.em_cartesian.EMComon_cofree}) and the Theorem-9.7
coalgebras $\FMeas(X)$
(\rocq{Icones.homs.em_cartesian.EMComon_FMeas}) are corollaries.
\end{theorem}

\begin{proof}
The naïve approach --- reduce everything to \emph{promoted points} via
\rocq{Icones.homs.em_seely_comonoid.d_bang_prom} --- cannot work for a
general $P = (A, a)$, because an image $a\,x$ is not promoted: the
defining computation of $d_A$ on $x!$ does not apply, and the laws fail
to reduce. We follow Melli\`es' \emph{structural} proof instead, a
retraction-and-lifting argument that never computes on points.

The first ingredient is that every coalgebra $P = (A, a)$ is a
\emph{retract} of its cofree coalgebra $\widetilde{\mathord{!}}\,UA$
inside $\mathrm{EM}(\mathord{!})$ (Melli\`es' Proposition 26): the
structure map $a : P \to \widetilde{\mathord{!}}\,UA$ is a coalgebra
morphism by the coassociativity of $P$, and the counit
$\varepsilon = \mathrm{der}_A$ retracts it
($\mathrm{der}_A \circ a = \mathrm{id}_A$ is the coalgebra counit law).
\rocq{Icones.homs.em_cartesian.diagram81} records the matching square
$(a \otimes a) \circ \mathrm{coalg\_d}\,P = d_{\mathord{!}A} \circ a$
on the nose --- not by point computation but by instantiating the
\emph{generally} valid \rocq{Icones.homs.em_cartesian.coalg_mor_d} at
the coalgebra morphism $a$, then folding through
\rocq{Icones.homs.em_cartesian.coalg_d_cofree}.

The second ingredient is that retracts of commutative comonoids are
commutative comonoids (Proposition 27). The four laws --- counit, left
counit, cocommutativity, coassociativity --- are transported along the
retraction by generic SMC cancellation
(\texttt{lcancel\_mono}) combined with the naturality of the associator,
braiding, and unitors against Diagram (85)
(\rocq{Icones.homs.em_cartesian.diagram81} above) and the cofree
comonoid laws on $\widetilde{\mathord{!}}\,UA$
(\autoref{lem:cbv-lc2}).

The hard step is showing the transported diagonal $\mathrm{coalg\_d}\,P$
is itself a coalgebra morphism into the product coalgebra $P \otimes P$
--- Melli\`es' Corollary 20, the one he flags as ``not so immediate''.
We package it as a generic \emph{lifting lemma}
\rocq{Icones.homs.em_cartesian.coalg_mor_lift} (Melli\`es' \S6.11
Proposition 20): if $i : X \to Y$ is a coalgebra morphism with a
retraction $r$ on the underlying carrier, and a candidate $f$ satisfies
$i \circ f$ is a coalgebra morphism, then $f$ is a coalgebra morphism.
Instantiated at $i = a \otimes a$ (a coalgebra morphism with retraction
$\mathrm{der} \otimes \mathrm{der}$) and $f = \mathrm{coalg\_d}\,P$, the
hypothesis $i \circ f = d_{\mathord{!}A} \circ a$ is exactly
\rocq{Icones.homs.em_cartesian.diagram81}, and a composite of two
coalgebra morphisms ($d_{\mathord{!}A}$ on the cofree side from
\autoref{lem:cbv-lc34}, and $a$); hence
$\mathrm{coalg\_d}\,P : P \to P \otimes P$ is a coalgebra morphism
(\rocq{Icones.homs.em_cartesian.coalg_d_is_mor_gen}). The easy half for
the counit is a composite of coalgebra morphisms
(\rocq{Icones.homs.em_cartesian.coalg_e_is_mor_gen}). Bundling the
four comonoid laws and the two coalgebra-morphism facts gives
\rocq{Icones.homs.em_cartesian.EMComon_all}.
\end{proof}

\begin{theorem}[$\mathrm{EM}(\mathord{!})$ is cartesian]%
\label{thm:cbv-cartesian}
\rocq{Icones.homs.em_cartesian.ICones_EM_cartesian}\rocqok
\uses{def:cbv-prod, lem:cbv-emcomon-all, def:cbv-em-cat}
The linear tensor $\otimes$ is the categorical product of the full
$\mathrm{EM}(\mathord{!})$ and the tensor unit $1$ is terminal. The
data --- projections \rocq{Icones.homs.em_cartesian.em_proj1_mor} /
\rocq{Icones.homs.em_cartesian.em_proj2_mor}, pairing
\rocq{Icones.homs.em_cartesian.em_pair} (no $\mathtt{EMComon}$
hypothesis needed, by \autoref{lem:cbv-emcomon-all}), the
$\beta$-laws \rocq{Icones.homs.em_cartesian.em_proj1_pair} /
\rocq{Icones.homs.em_cartesian.em_proj2_pair}, and the terminal
universal property \rocq{Icones.homs.em_cartesian.em_term_mor} /
\rocq{Icones.homs.em_cartesian.em_term_unique} --- is bundled in the
record \rocq{Icones.homs.em_cartesian.EM_Cartesian} with witness
\rocq{Icones.homs.em_cartesian.ICones_EM_cartesian} (axiom-free), whose
product/terminal field types reference the \emph{linear} $\otimes$, not
$\mathbin{\&}$.
\end{theorem}

\begin{proof}
Each coalgebra $P$ is comonoidal by \autoref{lem:cbv-emcomon-all}, so
Melli\`es' Corollary 17 applies: the projections, diagonal and pairing
out of $P$ are built from the comonoid data of $P$, $Q$. The
$\beta$-laws and terminal uniqueness follow from the comonoid laws of
\autoref{lem:cbv-lc2} as transported by
\autoref{lem:cbv-emcomon-all}.
\end{proof}

%%%%%%%%%%%%%%%%%%%%%%%%%%%%%%%%%%%%%%%%%%%%%%%%%%%%%%%%%%%%%%%%%%%%%%%%
\section{The cartesian \texorpdfstring{$\eta$}{eta}-rule on
\texorpdfstring{$\mathrm{EM}(\mathord{!})$}{EM(!)} (Fox 1976 /
Melli\`es Prop 28)}\label{sec:cbv-cart-eta}

The $\beta$-laws \rocq{Icones.homs.em_cartesian.em_proj1_pair} /
\rocq{Icones.homs.em_cartesian.em_proj2_pair} of
\autoref{thm:cbv-cartesian} establish the universal property of the
binary product \emph{out of} an arbitrary coalgebra $Z$. To complete the
cartesian structure we also need the $\eta$-rule
$\langle \pi_1, \pi_2 \rangle = \mathrm{id}$ at the linear level
$\ICones$. This is genuinely additional content --- Fox's 1976 theorem
that the EM category of a linear-exponential comonad is cartesian
under the linear tensor (Melli\`es Proposition 28 at the icones
level). The proof follows the same retract-and-lift recipe as
\autoref{lem:cbv-emcomon-all}: reduce to the cofree case, then transport
along the retraction.

\begin{lemma}[$\eta$-rule, cofree case]\label{lem:cbv-cart-eta-cofree}
\rocq{Icones.programs.infra.cbv_adjunction.em_pair_mor_proj_id_cofree}\rocqok
\uses{thm:cbv-cartesian, lem:tens-excl-charact}
For every $A, B : \ICones$, the pair-of-projections on the cofree
product equals the identity at the icones level:
\[
  \mathtt{em\_pair\_mor}\;(\pi_1)\;(\pi_2)
  \;=\;
  \mathrm{id}_{\mathord{!}A \otimes \mathord{!}B}.
\]
\end{lemma}

\begin{proof}
By \rocq{Icones.homs.em_seely_comonoid.tens_excl_charact}, it suffices
to check both sides agree on the promoted pure tensor $x! \otimes y!$
for $\cnorm{x} \leq 1$, $\cnorm{y} \leq 1$. On the cofree pair the
diagonal $\mathtt{coalg\_d}$ unfolds via
\rocq{Icones.homs.em_seely_comonoid.d_bang_prom} /
\rocq{Icones.homs.em_cartesian.m_bang_prom} /
\rocq{Icones.homs.bang.dig_prom} to
$(x! \otimes y!) \otimes (x! \otimes y!)$, and the projections
\rocq{Icones.homs.em_cartesian.em_proj1_mor} /
\rocq{Icones.homs.em_cartesian.em_proj2_mor} reduce on
$x! \otimes y!$ via the unitor + counit + \rocq{Icones.homs.em_seely_comonoid.e_bang_prom}
to $x!$ / $y!$ respectively.
\end{proof}

\begin{theorem}[The cartesian $\eta$-rule
$\langle \pi_1, \pi_2 \rangle = \mathrm{id}$]\label{thm:cbv-cart-eta}
\rocq{Icones.programs.infra.cbv_adjunction.em_pair_mor_proj_id}\rocqok
\uses{lem:cbv-cart-eta-cofree, lem:cbv-emcomon-all}
For every pair $(P, Q)$ of $\mathord{!}$-coalgebras, the
pair-of-projections of the binary product $P \otimes Q$ in
$\mathrm{EM}(\mathord{!})$ equals the identity on the underlying carrier:
\[
  \mathtt{em\_pair\_mor}\;(\pi_1)\;(\pi_2)
  \;=\;
  \mathrm{id}_{cP \otimes cQ}
  \quad : \quad
  \mathrm{Hom}_{\ICones}(cP \otimes cQ,\, cP \otimes cQ).
\]
Combined with the $\beta$-laws of \autoref{thm:cbv-cartesian}, this
completes the universal property of the EM(!) binary product
$(P \otimes Q,\, \pi_1,\, \pi_2,\, \langle \cdot, \cdot \rangle)$.
\end{theorem}

\begin{proof}
The structure maps
$(\mathtt{coalg\_str}\,P,\,\mathtt{coalg\_str}\,Q) : (P, Q) \to
(\widetilde{\mathord{!}}\,cP,\,\widetilde{\mathord{!}}\,cQ)$
are coalgebra morphisms (the coassociativity of $P$ and $Q$), and the
counit $\mathrm{der}_{cP} \otimes \mathrm{der}_{cQ}$ on the carrier
side retracts them ($\mathrm{der}_{cA} \circ a = \mathrm{id}_{A}$ is
the coalgebra counit). The tensored
$\mathtt{coalg\_str}\,P \otimes \mathtt{coalg\_str}\,Q$ on
$cP \otimes cQ$ is therefore a split mono with retraction
$\mathrm{der}_{cP} \otimes \mathrm{der}_{cQ}$. The $\eta$-rule on
$(P, Q)$ then follows from the cofree case
(\autoref{lem:cbv-cart-eta-cofree}) by naturality of $\pi_1, \pi_2$
along this retraction: $\langle \pi_1, \pi_2 \rangle$ at $(P, Q)$ is
the transport of $\langle \pi_1, \pi_2 \rangle$ at the cofree pair,
which equals the identity by \autoref{lem:cbv-cart-eta-cofree}, and the
retraction transports identities to identities.
\end{proof}

%%%%%%%%%%%%%%%%%%%%%%%%%%%%%%%%%%%%%%%%%%%%%%%%%%%%%%%%%%%%%%%%%%%%%%%%
\section{The monoidal adjunction and the CBV model}\label{sec:cbv-model}

The last step assembles the \emph{(lax symmetric) monoidal adjunction}
$U \dashv \widetilde{\mathord{!}}$ (Melli\`es' Proposition 29, via
Lack's lifting of the tensor to $\mathrm{EM}(\mathord{!})$), bundled with
the cartesian value category of \autoref{thm:cbv-cartesian}. Most of this
is assembly; the one block of new content is the lax-monoidal coherence
of $\widetilde{\mathord{!}}$ at the $m$-level, proved on promoted
points by the same $\texttt{tens\_excl\_charact}$ engine as
\autoref{ch:exponential}.

\begin{lemma}[Lax-monoidal coherence of the comparison
$m$]\label{lem:cbv-lax}
\rocq{Icones.programs.infra.cbv_adjunction.m_bang_assoc}\rocqok
\rocq{Icones.programs.infra.cbv_adjunction.m_bang_lunit}\rocqok
\rocq{Icones.programs.infra.cbv_adjunction.m_bang_runit}\rocqok
\rocq{Icones.programs.infra.cbv_adjunction.m_bang_braid}\rocqok
\uses{def:cbv-prod, lem:tens-excl-charact}
The comparison map $m$ (\autoref{def:cbv-prod}) satisfies the
associativity \rocq{Icones.programs.infra.cbv_adjunction.m_bang_assoc}, the unitors
\rocq{Icones.programs.infra.cbv_adjunction.m_bang_lunit} /
\rocq{Icones.programs.infra.cbv_adjunction.m_bang_runit} and the symmetry
\rocq{Icones.programs.infra.cbv_adjunction.m_bang_braid}, against the
$\alpha$ / $\lambda$ / $\rho$ / $\sigma$ of the tensor
(\autoref{thm:thm515}). Each is reduced by
\autoref{lem:tens-excl-charact} (and its ternary form) to agreement on
promoted pure tensors, then computed by
\rocq{Icones.homs.em_cartesian.m_bang_prom}.
\end{lemma}

\begin{theorem}[The CBV model on $\ICones$]\label{thm:cbv-model}
\rocq{Icones.programs.infra.cbv_adjunction.CBV_Model}\rocqok
\rocq{Icones.programs.infra.cbv_adjunction.ICones_CBV}\rocqok
\uses{thm:cbv-cofree-adj, thm:cbv-cartesian, lem:cbv-lax, thm:thm515}
The record \rocq{Icones.programs.infra.cbv_adjunction.CBV_Model} bundles the CBV /
LNL data on $\ICones$, mirroring the
\rocq{Icones.homs.seely.SeelyCategory} convention; its canonical witness
\rocq{Icones.programs.infra.cbv_adjunction.ICones_CBV} is \textbf{axiom-free}
(modulo the three classical \texttt{boolp} axioms, as for
\rocq{Icones.homs.seely.ICones_Seely}). It packages:
\begin{itemize}
  \item the symmetric monoidal closed structure
    $(\ICones, \otimes, 1)$ (\rocq{Icones.homs.smcc.ICones_SMCC},
    \autoref{thm:thm515});
  \item the Eilenberg--Moore category $\mathrm{EM}(\mathord{!})$
    (\rocq{Icones.homs.em_cat.ICones_EM}, \autoref{def:cbv-em-cat});
  \item the cartesian \emph{value} category --- the \emph{full}
    $\mathrm{EM}(\mathord{!})$ (\rocq{Icones.homs.em_cartesian.ICones_EM_cartesian},
    \autoref{thm:cbv-cartesian}) with its $\otimes$-product / $1$-terminal,
    unconditional comonoidality witness
    \rocq{Icones.homs.em_cartesian.EMComon_all}
    (\autoref{lem:cbv-emcomon-all}), and the cofree- /
    Theorem-9.7-generator corollaries
    \rocq{Icones.homs.em_cartesian.EMComon_cofree} /
    \rocq{Icones.homs.em_cartesian.EMComon_FMeas};
  \item the monoidal adjunction $U \dashv \widetilde{\mathord{!}}$
    over the \emph{full} $\mathrm{EM}(\mathord{!})$ (it needs no
    cartesianness): the functors, the (co)units, the hom-bijection with
    its round-trips and the triangle identities (\autoref{thm:cbv-cofree-adj});
  \item the monoidality witnesses: $U$ is strong (strict) monoidal --- the
    product / terminal carriers are the $\otimes$-tensor / $1$-unit, so the
    comparisons are identities --- and $\widetilde{\mathord{!}}$ is
    lax monoidal with comparison $m$
    (\rocq{Icones.programs.infra.cbv_adjunction.bang_m}) and unit
    (\rocq{Icones.programs.infra.cbv_adjunction.bang_e0}), its lax coherence
    (\autoref{lem:cbv-lax}) and the monoidal counit laws
    \rocq{Icones.programs.infra.cbv_adjunction.adj_counit_monoidal2} /
    \rocq{Icones.programs.infra.cbv_adjunction.adj_counit_monoidal0}.
\end{itemize}
\end{theorem}

\begin{proof}
Each field is one of the proved lemmas above; the witness
\rocq{Icones.programs.infra.cbv_adjunction.ICones_CBV} just assembles them. The
monoidal \emph{unit} laws \rocq{Icones.programs.infra.cbv_adjunction.adj_unit_monoidal2}
/ \rocq{Icones.programs.infra.cbv_adjunction.adj_unit_monoidal0} are recorded as
standalone facts (they are definitional --- the composite $m \circ (\eta
\otimes \eta)$ \emph{is} the product structure map), kept outside the
record to avoid the dependent-type plumbing. The cartesian fields now
hold on the \emph{full} $\mathrm{EM}(\mathord{!})$ unconditionally
(\autoref{lem:cbv-emcomon-all}), so $U \dashv \widetilde{\mathord{!}}$
is a genuine linear / non-linear adjunction with the (full) category of
$\mathord{!}$-coalgebras as the cartesian non-linear / value side.
\end{proof}

%%%%%%%%%%%%%%%%%%%%%%%%%%%%%%%%%%%%%%%%%%%%%%%%%%%%%%%%%%%%%%%%%%%%%%%%
\section{From the Moggi-CBV monad to the linhom-valued reading}%
\label{sec:cbv-calculus}

An earlier iteration of the formalisation accompanied
\autoref{thm:cbv-model} with a small first-order Moggi-CBV calculus
(intrinsically-typed Rocq inductives for unit, a measurable base
$\base\,X$, and binary products; values, computations
$\creturn V$ / $\clet x = M \mathbin{\mathtt{in}} N$ / $\csample V$),
interpreted into the CBV monad $T = \widetilde{\mathord{!}} \circ U$
of \autoref{thm:cbv-cofree-adj} with its Kleisli composition and
its $\mathtt{tunit\_eta}$ unit. That presentation has been retired:
the file \texttt{theories/programs/cbv.v} (which carried
$\mathtt{Tobj}$, $\mathtt{tunit\_eta}$, $\mathtt{kcomp}$, $\mathtt{kbind\_ext}$
and the three monad laws) is no longer part of the codebase, subsumed
by the linhom-valued, comonoid-primitive reading of
\autoref{sec:cbv-ppl}.

The two presentations agree on the underlying \emph{denotation} of
every closed term of the first-order calculus, but differ on the
\emph{shape} of the recursion. The Kleisli-monad reading interprets
every term as a coalgebra morphism into $T\,\sem{\tau}$ and recovers
identification rules ($\mathtt{let}$/$\mathtt{return}$, sample = the
integral) as monad-equational consequences. The linhom-valued reading
interprets every term as a \emph{linear} morphism
$U\sem{\Gamma} \to U\sem{\tau}$ (with no $T$ on the codomain) and
realises the sequencer $\mathtt{ne\_let}$, the pairing $\mathtt{ne\_pair}$
and the application $\mathtt{ne\_app}$ from the comonoid pair
$(\delta_\Gamma, \varepsilon_\Gamma)$ of every context coalgebra. The
$\mathord{!}$-induced monad on a measure type collapses the
$\csample V$ vs $\creturn V$ identification automatically: at
$\base\,X \mapsto \FMeas(X)$ the underlying iCone already carries
the Dirac inclusion, so a probabilistic sample from a Dirac and the
implicit return of its scalar value share the same linear morphism into
$\FMeas(X)$. The first-order calculus is therefore subsumed by the
higher-order linhom-valued interpreter of \autoref{sec:cbv-ppl}, and
no separate Moggi-CBV soundness theorem is recorded.

%%%%%%%%%%%%%%%%%%%%%%%%%%%%%%%%%%%%%%%%%%%%%%%%%%%%%%%%%%%%%%%%%%%%%%%%
\section{The boolean cone and the if-then-else combinator}%
\label{sec:cbv-bool}

\paragraph{This section is beyond the paper.}
To support the boolean elimination $\mathtt{ne\_if}$ in the higher-order
direct-style PPL (\autoref{sec:cbv-ppl}) we add a 2-point sub-probability
ICone, recognise it as paper \S4.4 / Theorem 4.24's coproduct
$\mathtt{cone\_one} \oplus \mathtt{cone\_one}$, package its universal
co-pairing as an icones-hom, and assemble an EM(!) Kleisli-level
$\mathtt{case\_em}$ combinator. The 2-point cone itself is a paper
construction (\S4.4); its concrete carrier
$\{\mathrm{nonneg}\,R\} \times \{\mathrm{nonneg}\,R\}$ through the HB
tower and the unit-ball-free icones-hom packaging are new infrastructure
in the spirit of the PPL.

\begin{definition}[The 2-point sub-probability cone]\label{def:cbv-bool-cone}
\rocq{Icones.programs.infra.bool_cone.bool_cone_car}\rocqok
\rocq{Icones.programs.infra.bool_cone.bool_dirac_true}\rocqok
\rocq{Icones.programs.infra.bool_cone.bool_dirac_false}\rocqok
\uses{def:icone}
The carrier $\mathtt{bool\_cone\_car}$ is a thin record
$\{\mathtt{bc\_t} : \{\mathrm{nonneg}\,R\};\; \mathtt{bc\_f} :
\{\mathrm{nonneg}\,R\}\}$ with pointwise operations and norm
$\cnorm{(p, q)} = p + q$; the unit ball is exactly the set of
sub-probability distributions on a 2-point space. The full HB tower
$\mathtt{isPrecone} \to \mathtt{isCone} \to \mathtt{isMCone} \to
\mathtt{isICone}$ is established, with the two Dirac points
$\mathtt{bool\_dirac\_true} = (1, 0)$ and
$\mathtt{bool\_dirac\_false} = (0, 1)$.
\end{definition}

\begin{remark}[Recognition as paper \S4.4 / Thm 4.24 coproduct]
\rocq{Icones.programs.infra.bool_cone.bool_case_true}\rocqok
\rocq{Icones.programs.infra.bool_cone.bool_case_false}\rocqok
The 2-point cone is the categorical coproduct
$\mathtt{cone\_one\_car} \oplus \mathtt{cone\_one\_car}$ of paper
\S4.4 / Theorem 4.24, with injections
$\mathtt{bool\_dirac\_true} / \mathtt{bool\_dirac\_false}$ playing the
role of $\mathtt{in\_t}(1) / \mathtt{in\_f}(1)$ and the universal
co-pairing $\mathtt{bool\_case}$ playing the role of $[a, b]$.
\end{remark}

\begin{definition}[The universal co-pairing
$\mathtt{bool\_case}$]\label{def:cbv-bool-case}
\rocq{Icones.programs.infra.bool_cone.bool_case}\rocqok
\uses{def:cbv-bool-cone}
For any iCone $A$ and any $a, b \in A$, set
$\mathtt{bool\_case}\;x\;a\;b\;=\;\mathtt{bc\_t}(x)\cdot a +
\mathtt{bc\_f}(x)\cdot b$. This is the unique
$\ICones$-morphism satisfying the coproduct universal property:
$\mathtt{bool\_case}\;\mathtt{bool\_dirac\_true}\;a\;b = a$ and
$\mathtt{bool\_case}\;\mathtt{bool\_dirac\_false}\;a\;b = b$.
\end{definition}

\begin{lemma}[Universal-property witnesses for
$\mathtt{bool\_case}$]\label{lem:cbv-bool-case-witnesses}
\rocq{Icones.programs.infra.bool_cone.bool_case_linear}\rocqok
\rocq{Icones.programs.infra.bool_cone.bool_case_omega_continuous}\rocqok
\rocq{Icones.programs.infra.bool_cone.bool_case_norm_le1}\rocqok
\rocq{Icones.programs.infra.bool_cone.bool_case_pres_path}\rocqok
\rocq{Icones.programs.infra.bool_cone.bool_case_pres_int}\rocqok
\uses{def:cbv-bool-case}
For $a, b \in A$ in the unit ball ($\cnorm{a} \leq 1$, $\cnorm{b} \leq 1$),
the map $x \mapsto \mathtt{bool\_case}\;x\;a\;b$ is linear, $\omega$-continuous,
of operator norm $\leq 1$, preserves measurable paths, and preserves
the Pettis integral against any finite measure. These are the
universal-property witnesses that upgrade the co-pairing to a full
$\ICones$ morphism (paper Def 4.10).
\end{lemma}

\begin{lemma}[Unit-ball-free variants and the test-measurability
generalisation]\label{lem:cbv-bool-case-gen}
\rocq{Icones.programs.infra.bool_cone.bool_case_omega_continuous_gen}\rocqok
\rocq{Icones.programs.infra.bool_cone.bool_case_norm_le_max}\rocqok
\rocq{Icones.programs.infra.bool_cone.bool_case_pres_path_gen}\rocqok
\rocq{Icones.programs.infra.bool_cone.bool_case_pres_int_gen}\rocqok
\rocq{Icones.mcones.mcone.test_meas_gen}\rocqok
\uses{lem:cbv-bool-case-witnesses}
The unit-ball restriction on $a, b$ can be dropped: $\omega$-continuity
of $\mathtt{bool\_case}$, path preservation, and integral preservation
hold unconditionally on $a, b$, with the operator norm bounded by
$\mathtt{Num.max}\;(\mathtt{Num.max}\;\cnorm{a}\;\cnorm{b})\;0$.
A general-cone-element variant of test measurability,
\rocq{Icones.mcones.mcone.test_meas_gen}, drops the unit-ball
hypothesis on the test scrutinee via scaling, exposing the test for
arbitrary cone elements.
\end{lemma}

\begin{definition}[Icones-hom packaging of $\mathtt{bool\_case}$]%
\label{def:cbv-bool-case-hom}
\rocq{Icones.programs.infra.bool_case_hom.bool_case_linhom}\rocqok
\rocq{Icones.programs.infra.bool_case_hom.bool_case_icones_hom}\rocqok
\rocq{Icones.programs.infra.bool_case_hom.bool_case_linhom_gen}\rocqok
\rocq{Icones.programs.infra.bool_case_hom.alpha_linhom}\rocqok
\rocq{Icones.programs.infra.bool_case_hom.beta_linhom}\rocqok
\uses{lem:cbv-bool-case-witnesses, lem:cbv-bool-case-gen}
The co-pairing packages as a $\mathtt{linhom\_car}$
(\rocq{Icones.programs.infra.bool_case_hom.bool_case_linhom}) and, via
\rocq{Icones.homs.seely.linhom_icones}, as a full
$\mathtt{icones\_hom}$
(\rocq{Icones.programs.infra.bool_case_hom.bool_case_icones_hom}), under the
unit-ball hypotheses on $a, b$. The unit-ball-free
$\mathtt{bool\_case\_linhom\_gen}$ drops these hypotheses. Although
$(a, b) \mapsto \mathtt{bool\_case}\;\cdot\;a\;b$ is not bilinear in
$(a, b)$, it decomposes as a sum
$\alpha(x, a) = \mathtt{bc\_t}(x)\cdot a$ plus
$\beta(x, b) = \mathtt{bc\_f}(x)\cdot b$ of two separately-bilinear
pieces; the linhom-level decomposition lemma is
\rocq{Icones.programs.infra.bool_case_hom.bool_case_linhom_gen_alpha_beta}.
\end{definition}

\begin{definition}[$\mathtt{if\_icones}$ --- the linhom-level
if-then-else combinator]\label{def:cbv-case-em}
\rocq{Icones.programs.ppl_cbv.if_icones}\rocqok
\uses{def:cbv-bool-case-hom, thm:cbv-cofree-adj, thm:cbv-cartesian}
The CBV interpreter of \autoref{sec:cbv-ppl} dispatches
$\mathtt{ne\_if}$ via the combinator
\rocq{Icones.programs.ppl_cbv.if_icones}. Given two branches as
$\mathtt{icones\_hom}$s $m, n : G \to A$ out of a common (carrier of a)
coalgebra context $G$ and into a common (carrier of a) target
coalgebra $A$, and a scrutinee
$b : G \to \widetilde{\mathord{!}}\,\mathtt{bool\_cone\_car}$, the
combinator builds the dispatch $\mathtt{if\_icones}\,m\,n\,b : G \to A$
by the recipe
\begin{enumerate}
  \item view the branches as norm-$\leq$-1 linhoms
    $G \multimap A$ via \rocq{Icones.homs.seely.icones_to_linhom};
  \item co-pair through
    \rocq{Icones.programs.infra.bool_case_hom.bool_case_linhom} into a
    norm-$\leq$-1 linhom
    $\mathtt{bool\_cone\_car} \multimap (G \multimap A)$;
  \item bridge back to an $\mathtt{icones\_hom}$ via
    \rocq{Icones.homs.seely.linhom_icones};
  \item SAFT-uncurry into $\mathtt{bool\_cone\_car} \otimes G \to A$;
  \item pre-compose with $(\mathrm{id}_G \otimes \mathrm{der})$ then a
    braid, taking $G \otimes \widetilde{\mathord{!}}\,\mathtt{bool\_cone\_car}$
    to $\mathtt{bool\_cone\_car} \otimes G$;
  \item pair $\mathrm{id}_G$ with the scrutinee through
    $\mathtt{em\_pair\_mor}$
    (\rocq{Icones.homs.em_cartesian.em_pair_mor}) to land at the
    target $G \to A$.
\end{enumerate}
The branches' norm-$\leq$-1 status comes for free from the
$\mathtt{icones\_hom}$ packaging (every icones-hom has operator norm
at most one), so the unit-ball hypotheses on
\rocq{Icones.programs.infra.bool_case_hom.bool_case_linhom} fire on the
nose. The earlier coalg-valued $\mathtt{case\_em}$ combinator of the
deleted \texttt{case\_em\_red.v} (Kleisli-reduction) is no longer
needed: in the linhom-valued reading the branches are bare linear
morphisms into the target's underlying iCone, and the dispatch closes
without a Kleisli round-trip.
\end{definition}

\paragraph{The \S9.7-style coalgebra structure on
$\mathtt{bool\_cone\_car}$.}
The interpretation $\sem{\mathtt{tbool}}$ of the PPL CBV interpreter
(\autoref{sec:cbv-ppl}) is \emph{not} the cofree fallback
$\widetilde{\mathord{!}}\,\mathtt{bool\_cone\_car}$: it is a non-cofree
$\mathord{!}$-coalgebra on $\mathtt{bool\_cone\_car}$ itself,
\rocq{Icones.programs.infra.bool_cone_coalg.bool_cone_coalg}, mirroring the
\S9.7-style \rocq{Icones.homs.coalgebra.FMeas_coalgebra} for the 2-point
case. This is the structural fix that gives the CBV
$\mathtt{ne\_bernoulli}$ / $\mathtt{ne\_if}$ pair its shared-sample
semantics.

\begin{definition}[The \S9.7-style coalgebra on
$\mathtt{bool\_cone\_car}$]\label{def:cbv-bool-cone-coalg}
\rocq{Icones.programs.infra.bool_cone_coalg.bool_coalg_str}\rocqok
\rocq{Icones.programs.infra.bool_cone_coalg.bool_cone_coalg}\rocqok
\uses{def:cbv-bool-cone, def:cbv-bool-case-hom}
The structure map
$\mathtt{bool\_coalg\_str} : \mathtt{bool\_cone\_car}\,Ar
\lin \mathord{!}\,\mathtt{bool\_cone\_car}\,Ar$
is the linhom-level universal co-pairing
\rocq{Icones.programs.infra.bool_case_hom.bool_case_icones_hom}
applied to the two promoted Dirac branches
$\mathrm{prom}(\delta_T)$ and $\mathrm{prom}(\delta_F)$, both unit-ball
by $\mathtt{prom\_ball}$ (\rocq{Icones.homs.bang.prom_ball}). On a
generic point $(p, q) \in \mathtt{bool\_cone\_car}$ it computes to
\[
  \mathtt{bool\_coalg\_str}\,(p, q)
  \;=\;
  p \cdot \mathrm{prom}(\delta_T)
  \;+\;
  q \cdot \mathrm{prom}(\delta_F)
  \quad\in\;
  \mathord{!}\,\mathtt{bool\_cone\_car}\,Ar,
\]
which is the finite-sum specialisation of the \S9.7 integral
$\mathtt{Coalg}_X(\mu) = \int \mathrm{prom}(\delta_x)\,d\mu(x)$
(\rocq{Icones.homs.coalgebra.FMeas_coalgebra}) when the underlying
measure space is the 2-point set $\{T, F\}$. The pair
$(\mathtt{bool\_cone\_car}\,Ar,\,\mathtt{bool\_coalg\_str})$ packages
as a \rocq{Icones.homs.coalgebra.Coalgebra} via
\rocq{Icones.programs.infra.bool_cone_coalg.bool_cone_coalg}, with sanity
witnesses
\rocq{Icones.programs.infra.bool_cone_coalg.bool_cone_coalg_obj} (the
carrier is $\mathtt{bool\_cone\_car}\,Ar$) and
\rocq{Icones.programs.infra.bool_cone_coalg.bool_cone_coalg_str_E} (the
structure map is $\mathtt{bool\_coalg\_str}$).
\end{definition}

\begin{lemma}[The two coalgebra laws]\label{lem:cbv-bool-coalg-laws}
\rocq{Icones.programs.infra.bool_cone_coalg.bool_coalg_counit}\rocqok
\rocq{Icones.programs.infra.bool_cone_coalg.bool_coalg_coassoc}\rocqok
\rocq{Icones.programs.infra.bool_cone_coalg.bool_cone_dispatch}\rocqok
\uses{def:cbv-bool-cone-coalg, def:cbv-cofree, thm:comonad}
The counit law
\rocq{Icones.programs.infra.bool_cone_coalg.bool_coalg_counit}
($\mathrm{der} \circ \mathtt{bool\_coalg\_str} = \mathrm{id}$) and the
coassociativity law
\rocq{Icones.programs.infra.bool_cone_coalg.bool_coalg_coassoc}
($\mathrm{dig} \circ \mathtt{bool\_coalg\_str} =
\mathord{!}(\mathtt{bool\_coalg\_str}) \circ \mathtt{bool\_coalg\_str}$)
are discharged via the universal-property \emph{dispatch} lemma
\rocq{Icones.programs.infra.bool_cone_coalg.bool_cone_dispatch}: by
linearity, two icones-homs $f, g : \mathtt{bool\_cone\_car} \to B$ are
equal iff they agree on the two basis points
$\delta_T$ and $\delta_F$, since
$x = \mathtt{bool\_case}\,x\,\delta_T\,\delta_F$ for every $x$
(\rocq{Icones.programs.infra.bool_cone_coalg.bool_cone_basis_expand}).
On the basis points, each side computes by the comonad identities
$\mathrm{der} \circ \mathrm{prom} = \mathrm{id}$
(\rocq{Icones.homs.bang.der_prom}) and
$\mathrm{dig} \circ \mathrm{prom} = \mathrm{prom} \circ \mathrm{prom}$
(\rocq{Icones.homs.bang.dig_prom}), combined with the basis identities
\rocq{Icones.programs.infra.bool_cone_coalg.bool_coalg_str_true} and
\rocq{Icones.programs.infra.bool_cone_coalg.bool_coalg_str_false}.
\end{lemma}

\paragraph{Shared-sample semantics for boolean draws.}
The transported commutative comonoid
$(\delta_{\mathtt{tbool}}, \varepsilon_{\mathtt{tbool}})$ that
\autoref{lem:cbv-emcomon-all} carries on the
$\mathord{!}$-coalgebra
\rocq{Icones.programs.infra.bool_cone_coalg.bool_cone_coalg} is the
\emph{diagonal pushforward}: under the CBV reading of
\autoref{sec:cbv-ppl}, the share-then-measure denotation of
\[
  \mathtt{let}\,x = \mathtt{Bernoulli}(p)\,\mathtt{in}\,(x, x)
\]
is $p \cdot (T, T) + (1 - p) \cdot (F, F)$, whereas the
independent-product denotation that the cofree fallback
$\widetilde{\mathord{!}}\,\mathtt{bool\_cone\_car}$ would induce is the
product measure
$p^2 \cdot (T, T) + p(1 - p) \cdot (T, F) +
(1 - p)p \cdot (F, T) + (1 - p)^2 \cdot (F, F)$. The diagonal-pushforward
rule is the one a CBV-style let on a freshly sampled boolean variable
should give, and is the boolean counterpart of the
$\mathtt{tbase}\,X$ shared-sample story carried by the
\rocq{Icones.homs.coalgebra.FMeas_coalgebra} of paper Theorem~9.7.

%%%%%%%%%%%%%%%%%%%%%%%%%%%%%%%%%%%%%%%%%%%%%%%%%%%%%%%%%%%%%%%%%%%%%%%%
\section{The CBV value-fixpoint at function types: $\mathtt{Yfix\_fun\_lin}$}%
\label{sec:cbv-yfix-fun}

The CBV interpreter resolves $\mathtt{ne\_fix}$ /
$\mathtt{ne\_fix\_mr}$ to a genuine Kleene value-fixpoint on the clean
linhom-cone $\mathrm{linhom}_\Ar(\Gamma,\,L)$ ($L$ the function-value
linhom-cone of the CBV reading, no $\widetilde{\mathord{!}} \circ U$
wrap on the codomain), built in
\texttt{theories/programs/infra/em\_fix.v}. The construction is
\emph{parametric} in the ambient coalgebra structure on $\Gamma$ --- it
only uses a diagonal $\delta : \Gamma \to \Gamma \otimes \Gamma$ on the
underlying icone --- so it applies uniformly to the cofree
$\widetilde{\mathord{!}}$-coalgebras at $\mathtt{tfun}$, the
\S9.7-shaped \rocq{Icones.homs.coalgebra.FMeas_coalgebra} at
$\mathtt{tbase}\,X$, the \rocq{Icones.programs.infra.bool_cone_coalg.bool_cone_coalg}
at $\mathtt{tbool}$, and the \rocq{Icones.homs.em_cartesian.EM_prod} at
$\mathtt{tprod}$. The arithmetic mirrors the SCones-side
\rocq{Icones.stable.fixpoint.Yfix} of \autoref{lem:cbn-recursion}: both
are a Kleene iteration on a unit-ball $\omega$-CPO, at the appropriate
category.

\subsection{The Kleene step $\mathtt{Phi\_fun}$}%
\label{ssec:cbv-phifun}

The step operator works at the linhom cone
$\mathrm{linhom}_\Ar(\Gamma,\,B)$, parametric in any icone-level
diagonal $\delta$ on $\Gamma$. There is no coalgebra structure required
on $B$: the codomain enters only through the body morphism
$M : \Gamma \otimes B \to B$.

\begin{definition}[$\mathtt{Phi\_fun}$ --- the CBV Kleene step]%
\label{def:cbv-phi-fun}
\rocq{Icones.programs.infra.em_fix.Phi_fun}\rocqok
\rocq{Icones.programs.infra.em_fix.Phi_fun_safe}\rocqok
\uses{thm:cbv-cartesian, def:cbv-prod}
For cones $\Gamma, B : \ICones$, a diagonal
$\delta : \Gamma \to \Gamma \otimes \Gamma$ and a body
$M : \Gamma \otimes B \to B$ (both packaged as $\ICones$ morphisms), the
step operator
\[
  \mathtt{Phi\_fun}\,M\,\mathit{prev}
  \;=\;
  M
    \;\circ\;
    \mathtt{tensor\_mor}\,(\mathrm{id}_\Gamma,\,\mathit{prev})
    \;\circ\;
    \delta
  \;\;:\;\;
  \Gamma \to B
\]
takes a previous approximant $\mathit{prev} : \Gamma \to B$ in the unit
ball of $\mathrm{linhom}_\Ar(\Gamma, B)$, duplicates the context through
$\delta$, runs the approximant on the second copy of $\Gamma$, and
applies the body. The total
\rocq{Icones.programs.infra.em_fix.Phi_fun} is defined on all of
$\mathrm{linhom}_\Ar(\Gamma, B)$ via $\mathtt{pselect}$: it agrees with the
safe variant
\rocq{Icones.programs.infra.em_fix.Phi_fun_safe} on the unit ball, and
returns $\mathtt{precone\_zero}$ off it (irrelevant for the Kleene
chain, which lives in the unit ball;
\rocq{Icones.programs.infra.em_fix.Phi_fun_unit}).
\end{definition}

\begin{lemma}[Ball preservation, monotonicity, and $\omega$-continuity
of $\mathtt{Phi\_fun}$]\label{lem:cbv-phi-fun-laws}
\rocq{Icones.programs.infra.em_fix.Phi_fun_ball}\rocqok
\rocq{Icones.programs.infra.em_fix.Phi_fun_incr}\rocqok
\rocq{Icones.programs.infra.em_fix.Phi_fun_cont}\rocqok
\uses{def:cbv-phi-fun}
On the unit ball,
$\mathtt{Phi\_fun}$ preserves the ball
(\rocq{Icones.programs.infra.em_fix.Phi_fun_ball}: if
$\cnorm{\mathit{prev}} \leq 1$ then
$\cnorm{\mathtt{Phi\_fun}\,M\,\mathit{prev}} \leq 1$, by composing the
three layer-norm bounds), is monotone
(\rocq{Icones.programs.infra.em_fix.Phi_fun_incr}: a precone-order
inequality $\mathit{prev}_1 \leq \mathit{prev}_2$ in the unit ball
transports through the three linear layers via
\rocq{Icones.cones.basic_lemmas.linear_increasing} and the layer-level
monotonicity witnesses, in particular
\rocq{Icones.programs.infra.em_fix.tensor_mor_R_lin_incr}), and is
$\omega$-continuous on the unit ball
(\rocq{Icones.programs.infra.em_fix.Phi_fun_cont}: it commutes with the
unit-ball supremum
\rocq{Icones.cones.cone.cone_sup_ball} of a unit-ball
chain of approximants, through the three layer-level continuity
witnesses
\rocq{Icones.programs.infra.em_fix.tensor_mor_omega_cont_R} on the
tensor layer,
\rocq{Icones.programs.infra.em_fix.linhom_pre_icones_sup} on the
$\delta$ pre-composition, and
\rocq{Icones.programs.infra.em_fix.linhom_post_icones_sup} on the $M$
post-composition).
\end{lemma}

\subsection{The Kleene supremum $\mathtt{Yfix\_fun\_lin}$}%
\label{ssec:cbv-yfix-sup}

The three laws of \autoref{lem:cbv-phi-fun-laws} make the unit-ball
$\omega$-CPO of $\mathrm{linhom}_\Ar(\Gamma, B)$ exactly the data the
generic linhom-level Kleene engine
\rocq{Icones.programs.infra.em_fix.linhom_lfp} needs.

\begin{definition}[The linhom-level Kleene chain
$\mathtt{kleene\_lin}$ and its supremum
$\mathtt{linhom\_lfp}$]\label{def:cbv-kleene-lin}
\rocq{Icones.programs.infra.em_fix.kleene_lin}\rocqok
\rocq{Icones.programs.infra.em_fix.linhom_lfp}\rocqok
\rocq{Icones.programs.infra.em_fix.linhom_lfp_norm_le1}\rocqok
\rocq{Icones.programs.infra.em_fix.linhom_lfp_fixpoint}\rocqok
\uses{lem:cbv-phi-fun-laws}
For any endo-operator $\Phi$ on $\mathrm{linhom}_\Ar(C, D)$ satisfying
the three laws of \autoref{lem:cbv-phi-fun-laws} (ball preservation,
monotonicity, $\omega$-continuity), the Kleene iterates
$\mathtt{kleene\_lin}\,\Phi\,n = \Phi^n(\mathtt{precone\_zero})$
form a chain in the unit ball
(\rocq{Icones.programs.infra.em_fix.kleene_lin_chain},
\rocq{Icones.programs.infra.em_fix.kleene_lin_ball}); their
unit-ball supremum is the least fixpoint
\rocq{Icones.programs.infra.em_fix.linhom_lfp}, with norm at most one
(\rocq{Icones.programs.infra.em_fix.linhom_lfp_norm_le1}) and the
fixpoint equation
$\Phi\,(\mathtt{linhom\_lfp}\,\Phi) = \mathtt{linhom\_lfp}\,\Phi$
(\rocq{Icones.programs.infra.em_fix.linhom_lfp_fixpoint}). This mirrors,
at the linhom cone, the SCones-side \rocq{Icones.stable.fixpoint.lfp} /
\rocq{Icones.stable.fixpoint.Yfix_fix} of the CBN chapter
(\autoref{lem:cbn-recursion}).
\end{definition}

\begin{definition}[The CBV value-fixpoint $\mathtt{Yfix\_fun\_lin}$]%
\label{def:cbv-yfix-fun-lin}
\rocq{Icones.programs.infra.em_fix.Yfix_fun_lin}\rocqok
\rocq{Icones.programs.infra.em_fix.Yfix_fun_lin_norm_le1}\rocqok
\rocq{Icones.programs.infra.em_fix.Yfix_fun_lin_fixpoint}\rocqok
\uses{def:cbv-phi-fun, def:cbv-kleene-lin}
For $\Gamma, B : \ICones$, a diagonal
$\delta : \Gamma \to \Gamma \otimes \Gamma$ and a body
$M : \Gamma \otimes B \to B$, the CBV value-fixpoint is
\[
  \mathtt{Yfix\_fun\_lin}\,\delta\,M
  \;=\;
  \sup_n (\mathtt{Phi\_fun}\,M)^n\,\mathtt{precone\_zero}
  \;\;:\;\;
  \Gamma \to B,
\]
the \rocq{Icones.programs.infra.em_fix.linhom_lfp} of
\rocq{Icones.programs.infra.em_fix.Phi_fun} at this $(\delta, M)$. Its
norm is at most one
(\rocq{Icones.programs.infra.em_fix.Yfix_fun_lin_norm_le1}) and it
satisfies the Kleene fixpoint equation
$\mathtt{Phi\_fun}\,M\,(\mathtt{Yfix\_fun\_lin}\,\delta\,M) =
\mathtt{Yfix\_fun\_lin}\,\delta\,M$
(\rocq{Icones.programs.infra.em_fix.Yfix_fun_lin_fixpoint}), proved by
specialising
\rocq{Icones.programs.infra.em_fix.linhom_lfp_fixpoint} to the three
laws of \autoref{lem:cbv-phi-fun-laws}.
\end{definition}

\subsection{Parametric in the coalgebra $B$}%
\label{ssec:cbv-yfix-parametric}

The Kleene engine of \autoref{def:cbv-yfix-fun-lin} works uniformly for
any iCone $B$ (and any diagonal $\delta$ on $\Gamma$). At the PPL CBV
interpreter level the diagonal is
$\delta = \mathtt{coalg\_d}\,(\mathtt{ctxD\_cbv}\,\Gamma)$ --- the
commutative-comonoid diagonal of \autoref{lem:cbv-emcomon-all} on the
underlying iCone of the context coalgebra --- and the codomain $B$ is
the underlying iCone of $\sem{\mathtt{tfun}\,A\,B}$. The same operator
applies regardless of \emph{which} of the four coalgebra shapes carries
$B$: the cofree $\widetilde{\mathord{!}}$-coalgebras
(\rocq{Icones.homs.em_cat.bang_cofree}) at $\mathtt{tfun}$, the
non-cofree \S9.7-shaped
\rocq{Icones.homs.coalgebra.FMeas_coalgebra} at $\mathtt{tbase}\,X$,
the non-cofree
\rocq{Icones.programs.infra.bool_cone_coalg.bool_cone_coalg} at
$\mathtt{tbool}$, and the \rocq{Icones.homs.em_cartesian.EM_prod}
product coalgebra at $\mathtt{tprod}$. This parametricity is what makes
the design symmetric across CBV types: the value-fixpoint never sees
the coalgebra structure on $B$, only its underlying iCone and the
ambient diagonal on $\Gamma$.

\subsection{Cross-reference to the CBN side}%
\label{ssec:cbv-yfix-cbn}

The arithmetic of $\mathtt{Yfix\_fun\_lin}$ mirrors the SCones-side
$\mathrm{Yfix}$ of \autoref{lem:cbn-recursion}: both build a Kleene
iteration on a unit-ball $\omega$-CPO from a body operator. The two
fixpoint operators land in different categories ($\ICones$-morphisms on
the linhom cone for CBV; $\mathbf{SCones}$-morphisms on the stable cone
for CBN) and feed into different surface clauses
(\rocq{Icones.programs.ppl_cbv.eD_cbv}'s $\mathtt{ne\_fix}$ for CBV;
\rocq{Icones.programs.ppl_cbn.eD_CBN_fix_E} for CBN), but share the
\emph{shape} of the recursion: take the supremum of
$\Phi^n(\mathtt{precone\_zero})$ on the unit ball.

\paragraph{Consequence for the recursive surface programs.}
The recursive examples \rocq{Icones.programs.examples.ex_loop},
\rocq{Icones.programs.examples.ex_geom} and
\rocq{Icones.programs.examples.ex_almost_loop} now have honest CBV
denotations through \rocq{Icones.programs.ppl_cbv.eD}: the
$\mathtt{ne\_fix}$ slot is the Kleene supremum of the body's iterates,
not the constant zero. Writing the closed-form mass identity for the
geometric program on the CBV side --- mirroring the axiom-free CBN
identity
\rocq{Icones.programs.ppl_cbn_geom.ex_geom_CBN_mass_one}
(\autoref{thm:cbn-ex-geom}) --- is a follow-up
(\autoref{sec:cbv-geom-mass}); the recursion infrastructure itself is
in place.

%%%%%%%%%%%%%%%%%%%%%%%%%%%%%%%%%%%%%%%%%%%%%%%%%%%%%%%%%%%%%%%%%%%%%%%%
\section{The higher-order direct-style PPL: linhom-valued interpretation}%
\label{sec:cbv-ppl}

The clean CBV interpreter \texttt{theories/programs/ppl\_cbv.v}
interprets the shared surface inductive
\rocq{Icones.programs.ppl.named_expr} as a \emph{linear} morphism in
$\ICones$, not as a coalgebra morphism. The shape of every clause is
fixed by the SMC primitives (tensor, braid, unitors, identity,
composition) and by the commutative-comonoid pair
$(\delta_\Gamma, \varepsilon_\Gamma) =
(\mathtt{coalg\_d}, \mathtt{coalg\_e})$ carried by every
$\mathord{!}$-coalgebra (Melli\`es Prop 28 / Cor 20,
\rocq{Icones.homs.em_cartesian.EMComon_all}). The cofree
$\widetilde{\mathord{!}}$ surfaces at exactly two clauses, both on the
function type: $\mathtt{ne\_lam}$ wraps via $\mathtt{adj\_psi}$, and
$\mathtt{ne\_app}$ unwraps via $\mathrm{der}$.

\paragraph{Why linhom-valued and not coalgebra-morphism-valued.}
Probabilistic programs do not in general commute with the
$\widetilde{\mathord{!}}$-coalgebra duplicability structure carried by
their source contexts; a probabilistic sample-then-pair has a different
denotation than a pair-then-sample. The minimal witness is the
let-then-pair fragment
\[
  \mathtt{let}\,x = \mathtt{sample}\,\mu\,\mathtt{in}\,(x, x)
  \;:\;
  \mathtt{tprod}\,\mathtt{tR}\,\mathtt{tR},
\]
whose share-then-measure denotation pushes $\mu$ along the diagonal
$r \mapsto (r, r)$ and whose measure-then-pair denotation is the
cartesian product measure $\mu \times \mu$; the two coincide iff $\mu$
is a Dirac. The Theorem-9.7 coalgebra structure on $\FMeas$ encodes
exactly the duplicability of shared samples, so a denotation that
\emph{is} a coalgebra morphism would force the share-then-measure
behaviour. The linhom-valued reading leaves the choice open: the
interpreter places no duplicability requirement on the denotation, and
the share-then-measure rule is the one a CBV-style let in fact gives.

\paragraph{Types and contexts.}
The type interpretation
\rocq{Icones.programs.ppl_cbv.tyD_cbv} sends every type to a coalgebra
of $\mathrm{EM}(\mathord{!})$:
\[
  \begin{array}{r@{\;\;=\;\;}l}
    \sem{\mathtt{tunit}}      & \EMterm \\
    \sem{\mathtt{tbase}\,X}   & \FMeas(X)
        \quad\text{(the Theorem-9.7 coalgebra)} \\
    \sem{\mathtt{tprod}\,t_1\,t_2}
                              & \EMprod\,\sem{t_1}\,\sem{t_2} \\
    \sem{\mathtt{tfun}\,A\,B} &
        \widetilde{\mathord{!}}\bigl(U\sem{A} \multimap U\sem{B}\bigr)
        \quad\text{(single $\widetilde{\mathord{!}}$, no $T$ on $B$)} \\
    \sem{\mathtt{tbool}}      &
        \mathtt{bool\_cone\_coalg}
        \quad\text{(the \S9.7-style 2-point coalgebra,
        \autoref{def:cbv-bool-cone-coalg}).}
  \end{array}
\]
Contexts are interpreted by
\rocq{Icones.programs.ppl_cbv.ctxD_cbv} right-to-left with the head on
the right, on the name-stripped skeleton
\rocq{Icones.programs.ppl.drop_names}~$\Gamma$; the variable witness
in $\mathtt{ne\_var}$ is projected to its De Bruijn index by
\rocq{Icones.programs.ppl.named_var_to_has_var} and then dispatched to
\rocq{Icones.programs.ppl_cbv.var_lookup_cbv}, a chain of
$\pi_1 / \pi_2$ projections over the tensor product context coalgebra.

The function-type clause is the only place where the source language
carries a single $\widetilde{\mathord{!}}$. There is no
$\widetilde{\mathord{!}} \circ U$ Moggi wrap on the codomain $B$: a
function value is a linear morphism
$U\sem{A} \multimap U\sem{B}$, and what is duplicable is the
\emph{value}, not the result of its application. This is the cleanest
reading the SMCC + EM(!) data of \autoref{thm:cbv-model} support.

\begin{remark}[$\mathtt{tbool}$ uses the \S9.7-style 2-point
coalgebra]\label{rem:cbv-tbool}
The clause $\sem{\mathtt{tbool}} = \mathtt{bool\_cone\_coalg}$ is
\emph{not} the cofree fallback
$\widetilde{\mathord{!}}\,\mathtt{bool\_cone\_car}$. It is the
non-cofree \S9.7-shaped $\mathord{!}$-coalgebra
\rocq{Icones.programs.infra.bool_cone_coalg.bool_cone_coalg} on the
carrier $\mathtt{bool\_cone\_car}\,Ar$ itself
(\autoref{def:cbv-bool-cone-coalg}), with structure map
$\mathtt{bool\_coalg\_str}\,(p, q) =
p \cdot \mathrm{prom}(\delta_T) + q \cdot \mathrm{prom}(\delta_F)$.
This is the boolean specialisation of paper Theorem~9.7's
$\mathtt{Coalg}_X(\mu) = \int \mathrm{prom}(\delta_x)\,d\mu(x)$ to the
2-point case --- which morally is
$\FMeas(\mathtt{bool\_meas\_space})$, hand-built here without
exposing $\mathtt{bool\_meas\_space} : \mathtt{measurableType}\,R$ as an
\rocq{Icones.mcones.ar.ar_obj} of $\mathbf{Ar}$ (a parameter-supplied
datum on $\mathbf{Ar}$). The transported commutative comonoid pair
$(\delta_{\mathtt{tbool}}, \varepsilon_{\mathtt{tbool}})$ on
$\mathtt{bool\_cone\_coalg}$ then recovers shared-sample semantics for
boolean draws (\autoref{def:cbv-bool-cone-coalg}).
\end{remark}

\paragraph{Term interpretation.}
The interpretation \rocq{Icones.programs.ppl_cbv.eD} assigns to every
$\Gamma \vdash M : \tau$ a linhom
\[
  \sem{M}
  \;:\;
  \mathrm{linhom}\bigl(U\sem{\mathtt{drop\_names}\,\Gamma},\,U\sem{\tau}\bigr),
\]
defined by structural recursion through the icones-hom-valued internal
helper \rocq{Icones.programs.ppl_cbv.eD_cbv}. Composing with
\rocq{Icones.homs.seely.icones_to_linhom} yields the public linhom; the
icones-hom variant lets every clause compose cleanly with
$\mathtt{tensor\_mor}$, $\mathtt{coalg\_d}$, $\mathtt{coalg\_e}$,
$\mathtt{tensor\_curry}$, $\mathtt{tensor\_uncurry}$. The
clause-by-clause shape is summarised below; in every line $\sem{M}$ is
an $\mathtt{icones\_hom}$ from $\sem{\Gamma}$'s carrier to
$\sem{\tau}$'s carrier.

\begin{itemize}
  \item $\mathtt{ne\_var}\,v$:
    \rocq{Icones.programs.ppl_cbv.var_lookup_cbv} of the variable
    witness $v$ --- a projection chain over the tensor product context.
  \item $\mathtt{ne\_tt}$: the comonoid counit
    $\varepsilon_\Gamma = \mathtt{em\_term\_mor}\,\sem{\Gamma}$
    (\rocq{Icones.homs.em_cartesian.em_term_mor}).
  \item $\mathtt{ne\_pair}\,M\,N$:
    $\mathtt{em\_pair\_mor}\,\sem{M}\,\sem{N}$
    (\rocq{Icones.homs.em_cartesian.em_pair_mor}); the diagonal
    $\delta_\Gamma = \mathtt{coalg\_d}$ is fired inside
    $\mathtt{em\_pair\_mor}$ to give each of the two sub-terms its own
    copy of the context.
  \item $\mathtt{ne\_fst}\,p$ / $\mathtt{ne\_snd}\,p$: pre-composition
    with \rocq{Icones.homs.em_cartesian.em_proj1_mor} /
    \rocq{Icones.homs.em_cartesian.em_proj2_mor}.
  \item $\mathtt{ne\_lam}\,x\,M$:
    $\mathtt{adj\_psi}\,(\mathtt{tensor\_curry}\,\sem{M})$. The body
    $\sem{M}$ is curried by the SMCC closure of $\ICones$
    (\autoref{thm:thm515}) into
    $\mathrm{linhom}(U\sem{\Gamma}, U\sem{A} \multimap U\sem{B})$, and
    $\mathtt{adj\_psi}$ of the cofree adjunction
    (\autoref{thm:cbv-cofree-adj}) at $\sem{\Gamma}$'s structure map
    promotes the result into
    $\widetilde{\mathord{!}}(U\sem{A} \multimap U\sem{B}) =
    \sem{\mathtt{tfun}\,A\,B}$. This is the first of the two
    $\widetilde{\mathord{!}}$ boundaries.
  \item $\mathtt{ne\_app}\,f\,e$:
    \[
      \sem{f \, e}
      \;=\;
      \mathtt{tensor\_uncurry}\,\mathrm{id}
        \;\circ\;
        (\mathrm{der}_{U\sem{A} \multimap U\sem{B}} \otimes \mathrm{id})
        \;\circ\;
        \mathtt{em\_pair\_mor}\,\sem{f}\,\sem{e}.
    \]
    The pair $\mathtt{em\_pair\_mor}\,\sem{f}\,\sem{e}$ supplies a copy
    of $\Gamma$ to each of $f$ and $e$ via $\delta_\Gamma$; the
    dereliction $\mathrm{der}$ unwraps the function value from its
    cofree $\widetilde{\mathord{!}}$ (the second $\widetilde{\mathord{!}}$
    boundary); $\mathtt{tensor\_uncurry}\,\mathrm{id}$ is the SMCC
    evaluation $(U\sem{A} \multimap U\sem{B}) \otimes U\sem{A} \to
    U\sem{B}$.
  \item $\mathtt{ne\_let}\,x\,M\,K$:
    $\sem{K} \circ \mathtt{em\_pair\_mor}\,\mathrm{id}_\Gamma\,\sem{M}$.
    The diagonal $\delta_\Gamma$ inside $\mathtt{em\_pair\_mor}$ pairs
    the context with the value of $M$, and the continuation $K$ is
    interpreted in the extended context. There is no Moggi-bind
    intermediate.
  \item $\mathtt{ne\_sample}\,\mu$, $\mathtt{ne\_real}\,r$,
    $\mathtt{ne\_true}$, $\mathtt{ne\_false}$,
    $\mathtt{ne\_bernoulli}\,p$: the constant $\mathtt{icones\_hom}$
    \rocq{Icones.programs.ppl.const_icones} out of $\sem{\Gamma}$ into,
    respectively, $\mu \in \FMeas(R_\mathrm{obj})$,
    $\delta_r \in \FMeas(R_\mathrm{obj})$, the Dirac points
    $\delta_T$ / $\delta_F$ in $\mathtt{bool\_cone\_car}$ (no
    $\mathrm{prom}$: the structure map of the
    \rocq{Icones.programs.infra.bool_cone_coalg.bool_cone_coalg}
    sends $\delta_T$ to $\mathrm{prom}(\delta_T)$ when the comonoid
    diagonal fires, but the constant lands in the carrier), and
    $\mathtt{bernoulli}\,p = p \cdot \delta_T + (1-p) \cdot \delta_F
    \in \mathtt{bool\_cone\_car}$. No coalgebra-morphism obligation
    arises: the constant linear morphism into a unit-ball element is
    automatically an icones-hom by
    \rocq{Icones.programs.ppl.const_icones}.
  \item $\mathtt{ne\_add}\,M\,N$ / $\mathtt{ne\_mul}\,M\,N$:
    \rocq{Icones.programs.ppl.add_lift} (resp.\
    \rocq{Icones.programs.ppl.mul_lift}) composed with
    $\mathtt{em\_pair\_mor}\,\sem{M}\,\sem{N}$. The arithmetic lifts
    are obtained from the FMeas lax-monoidal map
    (\autoref{thm:cbv-fmeas-lax}) post-composed with the
    $\FMeas$-functorial action of the measurable
    $\mathbb{R}\times\mathbb{R}\to\mathbb{R}$ operation.
  \item $\mathtt{ne\_score}\,f\,e$:
    $\mathtt{score\_lift}\,f \circ \sem{e}$, with
    \rocq{Icones.programs.ppl.score_lift} the density lift of
    \autoref{lem:ppl-score-tm-lift}.
  \item $\mathtt{ne\_if}\,b\,M\,N$:
    $\mathtt{if\_icones}\,\sem{M}\,\sem{N}\,\sem{b}$ of
    \autoref{def:cbv-case-em} --- the six steps use
    \rocq{Icones.programs.infra.bool_case_hom.bool_case_linhom},
    \rocq{Icones.homs.seely.linhom_icones},
    $\mathtt{tensor\_uncurry}$, a braid, and a final
    $\mathtt{em\_pair\_mor}$ with the scrutinee; with the §9.7-style
    $\mathtt{bool\_cone\_coalg}$ on $\sem{\mathtt{tbool}}$, the
    scrutinee lands directly in $\mathtt{bool\_cone\_car}$ (no
    $\mathrm{der}$ step required).
  \item $\mathtt{ne\_fix}\,s\,M$ (OCaml-style recursion at function
    types):
    \[
      \sem{\mathtt{ne\_fix}\,s\,M}
      \;=\;
      \mathtt{linhom\_icones}\,\bigl(
        \mathtt{Yfix\_fun\_lin}\;
          (\mathtt{coalg\_d}\,(\mathtt{ctxD\_cbv}\,\Gamma))\;
          (\sem{M})\bigr)\;
        \mathtt{Yfix\_fun\_lin\_norm\_le1},
    \]
    the CBV value-fixpoint of \autoref{def:cbv-yfix-fun-lin}, with the
    ambient context diagonal supplied by
    $\mathtt{coalg\_d}\,(\mathtt{ctxD\_cbv}\,\Gamma)$. The body
    $\sem{M}$ is interpreted in the context extended by the
    self-reference $s : \mathtt{tfun}\,A\,B$. The Kleene supremum is
    packaged back into an icones-hom via
    \rocq{Icones.homs.seely.linhom_icones}, using the unit-ball witness
    \rocq{Icones.programs.infra.em_fix.Yfix_fun_lin_norm_le1}.
  \item $\mathtt{ne\_fix\_mr}\,s\,t\,\_\,M$ (free-coalg-type
    generalisation, for mutual recursion / product-shaped recursive
    values): the same Kleene operator at the free-coalg-type case,
    \[
      \sem{\mathtt{ne\_fix\_mr}\,s\,t\,\_\,M}
      \;=\;
      \mathtt{linhom\_icones}\,\bigl(
        \mathtt{Yfix\_fun\_lin}\;
          (\mathtt{coalg\_d}\,(\mathtt{ctxD\_cbv}\,\Gamma))\;
          (\sem{M})\bigr)\;
        \mathtt{Yfix\_fun\_lin\_norm\_le1},
    \]
    parametric in $\Gamma$ and the coalgebra $t$, so the recipe is
    identical to $\mathtt{ne\_fix}$ above (see
    \autoref{ssec:cbv-yfix-parametric}).
\end{itemize}

\paragraph{The two $\widetilde{\mathord{!}}$ boundaries.}
The CBV interpreter never mentions
$\widetilde{\mathord{!}} \circ U$ in the codomain of any type, nor any
$\mathtt{Tobj}$ / $\mathtt{tunit\_eta}$ / Kleisli bind. The cofree
$\widetilde{\mathord{!}}$ appears in $\sem{\mathtt{tfun}\,A\,B}$ on the
outside, and the interpreter touches it exactly twice: $\mathtt{adj\_psi}$
at $\mathtt{ne\_lam}$ and $\mathrm{der}$ at $\mathtt{ne\_app}$. All
other compositional clauses ($\mathtt{ne\_let}$, $\mathtt{ne\_pair}$,
$\mathtt{ne\_fst}$ / $\mathtt{ne\_snd}$, $\mathtt{ne\_if}$,
$\mathtt{ne\_add}$ / $\mathtt{ne\_mul}$, $\mathtt{ne\_score}$) use only
the SMC primitives and the comonoid pair
$(\delta_\Gamma, \varepsilon_\Gamma)$. The diagonal $\delta_\Gamma$ is
what supplies each sub-expression of a compositional node with its own
copy of the context, so multi-use of a bound variable is free: it is a
passenger of the shared context, dispatched by the projection chain of
$\mathtt{var\_lookup\_cbv}$.

\paragraph{Arithmetic, score, and the FMeas lax-monoidal map.}
The interpretation of $\mathtt{ne\_add}$, $\mathtt{ne\_mul}$ and
$\mathtt{ne\_score}$ reuses the cross-cutting helpers
\rocq{Icones.programs.ppl.add_lift},
\rocq{Icones.programs.ppl.mul_lift} and
\rocq{Icones.programs.ppl.score_lift} of the surface module
\texttt{ppl.v}. These are $\ICones$-morphisms built from the FMeas lax
symmetric monoidal map of \autoref{thm:cbv-fmeas-lax} and the paper
\S6 follow-up
\rocq{Icones.homs.bilin.int_to_linhom_pres_path_in_cone}. Each binary
arithmetic clause composes the pair
$\mathtt{em\_pair\_mor}\,\sem{M}\,\sem{N}$ (whose diagonal is
$\delta_\Gamma$) with the lax-monoidal arithmetic lift; the score
clause $\mathtt{ne\_score}\,f\,e$ post-composes the density lift
$\mathtt{score\_lift}\,f$ with $\sem{e}$, producing the unit-cone
factor that scores the surrounding computation.

\begin{theorem}[The FMeas lax symmetric monoidal map at \texorpdfstring{$\ICones$}{ICones}]%
\label{thm:cbv-fmeas-lax}
\rocq{Icones.homs.fmeas_lax.fmeas_lax}\rocqok
\rocq{Icones.homs.fmeas_lax.fmeas_lax_E}\rocqok
\rocq{Icones.homs.fmeas_lax.fmeas_lax_dirac}\rocqok
\rocq{Icones.homs.bilin.int_to_linhom_pres_path_in_cone}\rocqok
\uses{thm:cbv-model}
For every $X, Y \in \mathbf{Ar}$ there is an $\ICones$ morphism
$\mathtt{fmeas\_lax}\,X\,Y :
   \FMeas(X) \otimes \FMeas(Y) \;\to\; \FMeas(X \times Y)$
that sends the pure tensor $\mu \otimes \nu$ to the cartesian
product measure $\mu \times \nu$ (\rocq{Icones.homs.fmeas_lax.fmeas_lax_E}),
and so sends Dirac point masses to the Dirac at the paired point,
$\mathtt{fmeas\_lax}\,(\delta_x \otimes \delta_y) = \delta_{(x,y)}$
(\rocq{Icones.homs.fmeas_lax.fmeas_lax_dirac}).
Construction. $\mathtt{fmeas\_lax}$ is obtained by
\rocq{Icones.homs.tensor_construct.tensor_uncurry} of the outer-linhom
$\FMeas(X) \multimap (\FMeas(Y) \multimap \FMeas(X \times Y))$
built by paper Theorem 6.1
(\rocq{Icones.homs.bilin.int_to_linhom}) twice: at the inner level
$\nu \mapsto \mathtt{fmeas\_lax\_pre}\,\delta_x\,\nu$ is a path
into $\FMeas(X \times Y)$ via the Dirac-pushforward path
$\mathtt{dirac\_lax}$, and at the outer level $x \mapsto (\nu \mapsto \cdots)$
is a path into the iCone $\FMeas(Y) \multimap \FMeas(X \times Y)$
of linhoms, which is the
\emph{path-preservation-in-the-cone-variable} step now discharged
as \rocq{Icones.homs.bilin.int_to_linhom_pres_path_in_cone} ---
the paper \S6 follow-up to Theorem 6.1 that the previous header of
\texttt{bilin.v} flagged as deferred.
\end{theorem}

\begin{proof}
\rocq{Icones.homs.fmeas_lax.fmeas_lax_E} reduces, on the pure
tensor, to the load-bearing Pettis/Tonelli identity
$\mathtt{fmeas\_lax\_pre\_iterated}$, expressing the pushforward
measure as the iterated icone-integral of $\mathtt{dirac\_lax}$.
The Dirac identity \rocq{Icones.homs.fmeas_lax.fmeas_lax_dirac}
then collapses to the function-level
$\mathtt{fmeas\_lax\_pre\_dirac}$, which is the elementary identity
$\delta_x \times_{\mathrm{meas}} \delta_y =
\delta_{(x,y)}$ along the measurable carrier cast
$\mathtt{ar\_prod\_cast}$. Construction of the outer linhom relies on
\rocq{Icones.homs.bilin.int_to_linhom_pres_path_in_cone}, which is
proved by packaging the cone-variable and the integrand into the
joint state $\mathbf{Ar}(Z) \times \mathbf{Ar}(Y)$ of
\rocq{Icones.icones.icone_integral.icone_integral_joint_measurable}
and specialising via the diagonal $(z, r) \mapsto (z, (z, r))$.
\end{proof}

\begin{lemma}[The term-level score lift]\label{lem:ppl-score-tm-lift}
\rocq{Icones.programs.ppl.score_lift}\rocqok
\rocq{Icones.programs.ppl.score_lift_dirac}\rocqok
\uses{thm:cbv-fmeas-lax}
For every measurable $f : \mathbb{R} \to \mathbb{R}$ with
$0 \le f(r) \le 1$ for all $r$, there is an $\ICones$ morphism
$\mathtt{score\_lift}\,f :
   \FMeas(R_{\mathrm{obj}}) \;\to\; \mathtt{cone\_one\_car}$
of operator norm at most one, satisfying the Dirac identity
$\mathtt{score\_lift}\,f\,(\delta_{\mathtt{R\_to\_carrier}\,r})
= f(r) \cdot \mathtt{one1}$
(\rocq{Icones.programs.ppl.score_lift_dirac}).
\end{lemma}

\begin{proof}
The path $r \mapsto f(r) \cdot \mathtt{one1}$ is a measurable path
into $\mathtt{cone\_one\_car}$ of norm at most one (pointwise
$f(r) \le 1$), so by paper Theorem 6.1
(\rocq{Icones.homs.bilin.int_to_linhom}) it lifts to an
$\FMeas(R_{\mathrm{obj}}) \multimap \mathtt{cone\_one\_car}$ linhom
$\mathtt{int\_to\_linhom}\,\mathtt{score\_path}$; the
$\ICones$-morphism packaging $\mathtt{linhom\_icones}$ then yields
$\mathtt{score\_lift}$. The Dirac identity follows by integration
against a Dirac evaluating the path at the carrier point.
\end{proof}

\paragraph{The shared surface programs.}
The eleven surface PPL programs in
\texttt{theories/programs/examples.v} type-check under
\rocq{Icones.programs.ppl_cbv.eD}: the recursive examples
\rocq{Icones.programs.examples.ex_loop},
\rocq{Icones.programs.examples.ex_geom} and
\rocq{Icones.programs.examples.ex_almost_loop} now resolve, at the
$\mathtt{ne\_fix}$ slot, to the genuine Kleene fixpoint
$\mathtt{Yfix\_fun\_lin}$ of \autoref{def:cbv-yfix-fun-lin}, not to a
$\mathtt{precone\_zero}$ placeholder. The QBS-style examples ---
\rocq{Icones.programs.examples.ex_random_constant} (a distribution
over a function space, with source type
$\mathtt{tfun}\,\mathtt{tR}\,\mathtt{tR}$),
\rocq{Icones.programs.examples.ex_random_linear} (a distribution over
$\lambda x.\,m\cdot x + b$ for $m, b \sim \mu$, exercising
$\mathtt{ne\_add}$ and $\mathtt{ne\_mul}$), and
\rocq{Icones.programs.examples.ex_bayes_linear} (a prior/score/observe
shape exercising $\mathtt{ne\_score}$) --- close fully on
$\mathtt{ne\_let}$, $\mathtt{ne\_sample}$, $\mathtt{ne\_lam}$,
$\mathtt{ne\_app}$, $\mathtt{ne\_add}$, $\mathtt{ne\_mul}$ and
$\mathtt{ne\_score}$. The boolean sanity examples
\rocq{Icones.programs.examples.ex_true},
\rocq{Icones.programs.examples.ex_false},
\rocq{Icones.programs.examples.ex_fair_coin}, and
\rocq{Icones.programs.examples.ex_if_demo} exercise the
$\mathtt{if\_icones}$ dispatch of \autoref{def:cbv-case-em}. Structural
$\mathtt{\_denot\_E}$ reduction lemmas of the kind previously hosted
in \texttt{examples.v} against the Moggi-monad interpretation have
been retired; the analogous reductions in the linhom-valued reading
unfold to the clause specifications above without any monad-equational
manipulation, and are recovered by definitional unfolding of
\rocq{Icones.programs.ppl_cbv.eD_cbv}.

\paragraph{Named-variable surface syntax.}
The named surface is the only syntax of
\texttt{theories/programs/ppl.v}: no separate De Bruijn calculus, no
intermediate translation. Variable lookup $\#\,\texttt{"x"}$ uses
canonical structures
(\rocq{Icones.programs.ppl.find_nv}, with
\rocq{Icones.programs.ppl.found_nctx} /
\rocq{Icones.programs.ppl.recurse_nctx} carrying the search state and
\texttt{ne\_var'} delivering the resolved
$\mathtt{ne\_var}$ at a $\mathtt{find\_nv}$ witness) with
\mathcompA{}'s \texttt{infer} typeclass on \texttt{String.eqb}, so
Coq's elaborator infers the context slot, type, and named-variable
witness simultaneously. The custom entry \texttt{ppl\_named} supports
\texttt{let "x" := M in N},
\texttt{\textbackslash "x" ::: A => M},
\texttt{fix "f" ::: tfun A B in M} (desugaring to
\rocq{Icones.programs.ppl.ne_fix}, the genuine Kleene value-fixpoint
of \autoref{def:cbv-yfix-fun-lin}),
\texttt{Sample (mu, Hmu)},
\texttt{Score \{ f, Hf\_meas, Hf\_ge0, Hf\_le1 \} e},
\texttt{True} / \texttt{False},
\texttt{Bernoulli \{ p, Hp\_ge0, Hp\_le1 \}},
\texttt{if e then M else N} (dispatched via
\autoref{def:cbv-case-em}),
\texttt{\# "x"}, \texttt{M @ N}, \texttt{M + N}, \texttt{M * N},
\texttt{(e1, e2)}, \texttt{fst e}, \texttt{snd e}, \texttt{()},
\texttt{[|r|]}, and \texttt{\{x\}}-escape. There is no
\texttt{Ret} surface form. All eleven programs in
\texttt{theories/programs/examples.v} are written exclusively in this
direct-style surface notation.

\paragraph{Out-of-scope items.}
Adequacy, normalization, and full-abstraction (traditional metatheory,
requiring more than a denotational model). Coproducts beyond the
2-point case of \autoref{sec:cbv-bool} (no general sum type in the
cones model). A closed-form CBV-side mass identity for the recursive
surface programs (\autoref{sec:cbv-geom-mass}) --- the recursion
infrastructure of \autoref{def:cbv-yfix-fun-lin} is now in place, but
the specific mass identity mirroring the CBN-side
\rocq{Icones.programs.ppl_cbn_geom.ex_geom_CBN_mass_one}
(\autoref{thm:cbn-ex-geom}) has not yet been written. A Moggi
metalanguage with an explicit value/computation split (orthogonal to
the linhom-valued framing adopted here).

%%%%%%%%%%%%%%%%%%%%%%%%%%%%%%%%%%%%%%%%%%%%%%%%%%%%%%%%%%%%%%%%%%%%%%%%
\section{Mass identity for \texorpdfstring{$\mathtt{ex\_geom}$}{ex\_geom} on the CBV side}%
\label{sec:cbv-geom-mass}

A CBV-side closed-form mass identity for the geometric program
\rocq{Icones.programs.examples.ex_geom} would parallel the
CBN-side $\rocq{Icones.programs.ppl_cbn_geom.ex_geom_CBN_mass_one}$
of \autoref{thm:cbn-ex-geom}: the Kleene cascade on the unit-ball
$\omega$-CPO of \autoref{def:cbv-yfix-fun-lin} yields per-iterate mass
$1 - (1/2)^n$ (Bernoulli at $1/2$ pulls $\delta_0$ with weight $1/2$
and recurses on the remaining $1/2$ through the unary linear
pushforward of the ELSE branch), with limit $1$ by monotone
convergence in $\overline{\mathbb{R}}$. With $\mathtt{ne\_fix}$ now
routed through $\mathtt{Yfix\_fun\_lin}$ in
\rocq{Icones.programs.ppl_cbv.eD_cbv}, the CBV-side identity is
provable in principle on the current infrastructure --- the recursion
shape is in place and the per-iterate cascade and mass-convergence
lemmas of \autoref{lem:cbn-kleene-cascade} /
\autoref{lem:cbn-kleene-cvg} have direct CBV analogues at the linhom
cone. The specific CBV proof has not yet been written and is a
follow-up. The axiom-free CBN-side identity
\rocq{Icones.programs.ppl_cbn_geom.ex_geom_CBN_mass_one} of
\autoref{ch:cbn} remains the current closed record for the geometric
program in the cones model.

%%%%%%%%%%%%%%%%%%%%%%%%%%%%%%%%%%%%%%%%%%%%%%%%%%%%%%%%%%%%%%%%%%%%%%%%
\section{Status}\label{sec:cbv-status}

\paragraph{Delivered and axiom-free.}
The CBV \emph{model} and its supporting infrastructure are axiom-clean
(only the three classical \texttt{boolp} axioms): the Eilenberg--Moore
category and cofree adjunction (\autoref{sec:cbv-em},
\rocq{Icones.homs.em_cat.ICones_EM}); the Seely comonoid with LC2--LC4
(\autoref{sec:cbv-comonoid}); the \emph{full} cartesian
$\mathrm{EM}(\mathord{!})$ carried by the linear $\otimes$
(\autoref{thm:cbv-cartesian},
\rocq{Icones.homs.em_cartesian.ICones_EM_cartesian}) with the
unconditional comonoidality witness
\rocq{Icones.homs.em_cartesian.EMComon_all}
(\autoref{lem:cbv-emcomon-all}); the cartesian $\eta$-rule
\rocq{Icones.programs.infra.cbv_adjunction.em_pair_mor_proj_id}
(\autoref{thm:cbv-cart-eta}, Fox 1976 / Melli\`es Prop 28 at the
icones level); the LNL monoidal adjunction bundled as
\rocq{Icones.programs.infra.cbv_adjunction.ICones_CBV}
(\autoref{thm:cbv-model}); the FMeas lax symmetric monoidal map
\rocq{Icones.homs.fmeas_lax.fmeas_lax} at the $\ICones$ level together
with its Dirac identity \rocq{Icones.homs.fmeas_lax.fmeas_lax_dirac}
and the discharged paper-\S6 follow-up
\rocq{Icones.homs.bilin.int_to_linhom_pres_path_in_cone}
(\autoref{thm:cbv-fmeas-lax}); the term-level score lift
\rocq{Icones.programs.ppl.score_lift} and its Dirac identity
\rocq{Icones.programs.ppl.score_lift_dirac}
(\autoref{lem:ppl-score-tm-lift}); the boolean cascade
(\autoref{sec:cbv-bool}) ---
\rocq{Icones.programs.infra.bool_cone.bool_cone_car} as the paper
\S4.4 / Thm 4.24 coproduct
$\mathtt{cone\_one}\oplus\mathtt{cone\_one}$, the icones-hom packaging
\rocq{Icones.programs.infra.bool_case_hom.bool_case_icones_hom} (with
unit-ball-free variant
\rocq{Icones.programs.infra.bool_case_hom.bool_case_linhom_gen}), the
linhom-level if-then-else combinator
\rocq{Icones.programs.ppl_cbv.if_icones} used by $\mathtt{ne\_if}$
(\autoref{def:cbv-case-em}), and the \S9.7-style 2-point
$\mathord{!}$-coalgebra
\rocq{Icones.programs.infra.bool_cone_coalg.bool_cone_coalg}
(\autoref{def:cbv-bool-cone-coalg}) used as the carrier of
$\sem{\mathtt{tbool}}$; the CBV value-fixpoint
\rocq{Icones.programs.infra.em_fix.Yfix_fun_lin}
(\autoref{def:cbv-yfix-fun-lin}) on the clean linhom-cone, built from
the Kleene step
\rocq{Icones.programs.infra.em_fix.Phi_fun} of
\autoref{def:cbv-phi-fun} via the generic linhom-level engine
\rocq{Icones.programs.infra.em_fix.linhom_lfp}
(\autoref{def:cbv-kleene-lin}), parametric in the ambient context
diagonal and in the codomain coalgebra shape
(\autoref{ssec:cbv-yfix-parametric}); the linhom-valued, comonoid-primitive
CBV interpreter \rocq{Icones.programs.ppl_cbv.eD} of
\autoref{sec:cbv-ppl}, with its type interpretation
\rocq{Icones.programs.ppl_cbv.tyD_cbv}, context interpretation
\rocq{Icones.programs.ppl_cbv.ctxD_cbv}, and variable lookup
\rocq{Icones.programs.ppl_cbv.var_lookup_cbv}; and the shared surface
syntax of \texttt{theories/programs/ppl.v} consumed by both this CBV
interpreter and the CBN interpreter of \autoref{ch:cbn}, with the
eleven surface programs of \texttt{theories/programs/examples.v}
(\rocq{Icones.programs.examples.ex_random_constant},
\rocq{Icones.programs.examples.ex_random_linear},
\rocq{Icones.programs.examples.ex_bayes_linear},
\rocq{Icones.programs.examples.ex_loop},
\rocq{Icones.programs.examples.ex_geom},
\rocq{Icones.programs.examples.ex_almost_loop}, plus the four boolean
sanity checks).

\paragraph{Follow-ups.}
The recursion infrastructure of \autoref{sec:cbv-yfix-fun} is now in
place, so the only CBV-side residue listed in earlier drafts of this
status section --- the $\mathtt{ne\_fix}$ / $\mathtt{ne\_fix\_mr}$
placeholders and the cofree-based $\mathtt{tbool}$ clause --- has
landed (the former through
\rocq{Icones.programs.infra.em_fix.Yfix_fun_lin}, the latter through
\rocq{Icones.programs.infra.bool_cone_coalg.bool_cone_coalg}). What
remains as follow-up is a specific CBV-side mass identity for the
recursive surface programs, mirroring the axiom-free CBN-side
\rocq{Icones.programs.ppl_cbn_geom.ex_geom_CBN_mass_one}
(\autoref{thm:cbn-ex-geom}); the per-iterate cascade and
mass-convergence shape are already in place at the linhom cone via the
ω-continuity witness
\rocq{Icones.programs.infra.em_fix.Phi_fun_cont} and the unit-ball
supremum \rocq{Icones.cones.cone.cone_sup_ball}, but writing the
specific proof for \rocq{Icones.programs.examples.ex_geom} on the new
\rocq{Icones.programs.infra.em_fix.Yfix_fun_lin} has not yet been
attempted (\autoref{sec:cbv-geom-mass}).

\paragraph{What is not interpreted.}
$\mathrm{EM}(\mathord{!})$ is not cartesian closed and is not expected
to be --- a structural fact about EM categories of linear-exponential
comonads, not a missing diagram chase. The linhom-valued reading of
\autoref{sec:cbv-ppl} is what the SMCC + EM(!) data of
\autoref{thm:cbv-model} naturally supports, and that holds axiom-free.
The genuinely out-of-scope items are: adequacy, normalization, and
full-abstraction (traditional metatheory, requiring more than a
denotational model); coproducts beyond the 2-point boolean case of
\autoref{sec:cbv-bool} (no general sum type in the cones model); a
$\mathtt{qbs\_normalize}$ / Bayesian-posterior pass (a different
observable than the unnormalised prior/score/observe shape of
\rocq{Icones.programs.examples.ex_bayes_linear}); and a Moggi
metalanguage with an explicit value/computation split (orthogonal to
the linhom-valued framing adopted here).
